%%%%%%%%%%%%%%%%%%%%%%%%%%%%%%%%%%%%%%%%%%%%%%%%%%%%%%%%%%%%%%%%%%%%%%
% Lớp 12 | Phương trình mặt phẳng, đường thẳng, mặt cầu | Phương trình mặt phẳng
% Dạng 2H5?1-5: Khoảng cách điểm tới mặt phẳng
% Số câu: 58  (sinh tự động bởi tach_cau.py; từ nay thầy sửa trực tiếp file này)
%%%%%%%%%%%%%%%%%%%%%%%%%%%%%%%%%%%%%%%%%%%%%%%%%%%%%%%%%%%%%%%%%%%%%%

%%%%%%%%%%%%%%%% TRẮC NGHIỆM 4 PHƯƠNG ÁN %%%%%%%%%%%%%%%%

% nguồn: Phan1_Lop_12_Hinh_hoc_Chuong_5
\begin{ex}%[2H5N1-5]%[ma:62cd9dc0]
	Trong không gian với hệ tọa độ $Oxyz$,  khoảng cách từ điểm $M(9;7;8)$đến mặt phẳng $(P)\colon ax+by+cz+d=0$ bằng
	\choice 
		{$\dfrac{|7a+8b+9c+d|}{\sqrt{a^2+b^2+c^2}}$} 
		{\True $\dfrac{|9a+7b+8c+d|}{\sqrt{a^2+b^2+c^2}}$} 
		{$\dfrac{|7a+8b+9c+d|}{\sqrt{7^2+8^2+9^2}}$} 
		{$\dfrac{|9a+7b+8c+d|}{\sqrt{9^2+7^2+8^2}}$} 
	 \loigiai{Ta có $\mathrm{d}\big(M,(P)\big)=\dfrac{|9a+7b+8c+d|}{\sqrt{a^2+b^2+c^2}}$.}
\end{ex}

% nguồn: Phan1_Lop_12_Hinh_hoc_Chuong_5
\begin{ex}%[2H5N1-5]%[ma:3d27724a]
	Trong không gian $Oxyz$, cho điểm $A(1; 2; 2)$ và mặt phẳng  $(\alpha)\colon x+2y-2z-4=0$. Khoảng cách từ điểm $A$ đến mặt phẳng $(\alpha)$ bằng
	\choice 
		 {$3$}
		 {\True $1$}
		{$\dfrac{13}{3}$} 
		 {$\dfrac{1}{3}$}
	\loigiai{ 
		Ta có  $\mathrm{d}\left(A,(\alpha)\right)=\dfrac{|1+2\cdot2-2\cdot2-4|}{\sqrt{1^2+2^2+(-2)^2}}=1$.
	} 
\end{ex}

% nguồn: Phan1_Lop_12_Hinh_hoc_Chuong_5
\begin{ex}%[2H5N1-5]%[ma:c341ad54]
	Trong không gian $Oxyz$, khoảng cách từ $M(1;2;-3)$ đến  $(P)\colon x+2y+2z-10=0$ là
	\choice 
		 {$3$}
		 {$\dfrac{2}{3}$}
		{$\dfrac{4}{3}$} 
		{\True  $\dfrac{11}{3}$}
	\loigiai{ 
	Ta  có $\mathrm{d}\left(M,(P)\right)=\dfrac{|1+4-6-10|}{\sqrt{1+4+4}}=\dfrac{11}{3}$.}
\end{ex}

% nguồn: Phan1_Lop_12_Hinh_hoc_Chuong_5
\begin{ex}%[2H5N1-5]%[ma:ed07949f]
	Trong không gian $Oxyz$, cho $M(3;-1;-2)$ và mặt phẳng $(\alpha)\colon 3x-y+2z+4=0$. Khoảng cách từ điểm $M$ đến mặt phẳng $(\alpha)$ bằng 
	 \choice
		{$\dfrac{5\sqrt{13}}{7}$} 
		 {$\dfrac{5\sqrt{15}}{7}$}
		{\True  $\dfrac{5\sqrt{14}}{7}$}
		{$\dfrac{5\sqrt{17}}{7}$} 
	 \loigiai{
		$\mathrm{\,d}\left(M;(\alpha)\right)=\dfrac{|3\cdot3+1-4+4|}{\sqrt{14}}=\dfrac{5\sqrt{14}}{7}$. 
	 }
\end{ex}

% nguồn: Phan1_Lop_12_Hinh_hoc_Chuong_5
\begin{ex}%[2H5N1-5]%[ma:30c81f98]
	Trong không gian $Oxyz$, khoảng cách  từ điểm $A(1;1;0)$ đến mặt phẳng $(\alpha) \colon 2x+2y+z-1=0$ bằng
	\choice 
		{\True $1$} 
		 {$\dfrac{4}{3}$}
		{$\dfrac{3\sqrt{2}}{2}$} 
		 {$2\sqrt{2}$}
	 \loigiai{
		Khoảng cách từ điểm $A$ đến mặt phẳng $(\alpha)$~là 
		$$\mathrm{d}\left(A,(\alpha)\right)=\dfrac{\left|2\cdot1+2\cdot1+0-1\right|}{\sqrt{2^2+2^2+1^2}}=1$$ 
	 }
\end{ex}

% nguồn: Phan1_Lop_12_Hinh_hoc_Chuong_5
\begin{ex}%[2H5N1-5]%[ma:88371775]
	Trong không gian $Oxyz$, khoảng cách từ  điểm $K(2; 0;-1)$ đến mặt phẳng $(P)\colon 2 x-y+z-1=0$ bằng
	 \choice
		{\True  $\dfrac{2}{\sqrt{6}}$}
		{$\dfrac{\sqrt{6}}{2}$} 
		 {$\dfrac{1}{\sqrt{6}}$}
		{$\dfrac{\sqrt{6}}{6}$} 
	 \loigiai{
		$\mathrm{d}(K,(P))=\dfrac{|2\cdot 2-0-1-1|}{\sqrt{2^2+(-1)^2+1^2}}=\dfrac{2}{\sqrt{6}}$. 
	 }
\end{ex}

% nguồn: Phan1_Lop_12_Hinh_hoc_Chuong_5
\begin{ex}%[2H5N1-5]%[ma:71cfebae]
	Trong không gian $Oxyz$, khoảng cách từ điểm $A(-2;0;2)$ đến mặt phẳng $(P)$ có phương trình $2x - y + 2z - 6 = 0$ là 
	\choice 
		 {$6$}
		{$9$} 
		{$3$} 
		{\True  $2$}
	\loigiai{ 
		Khoảng cách  từ điểm $A$ đến mặt phẳng $(P)$ là
	\[\mathrm{d}\left(A,(P)\right)=\dfrac{|2 \cdot (-2) -  0 +2 \cdot 2-6| }{\sqrt{2^2+(-1)^2+2^2}}=2.\]}
\end{ex}

% nguồn: Phan1_Lop_12_Hinh_hoc_Chuong_5
\begin{ex}%[2H5H1-5]%[ma:b2c01ebf]
	Trong  không gian $Oxyz$, khoảng cách giữa hai mặt phẳng $(P)\colon x+2y+2z-10=0$ và $(Q)\colon x+2y+2z-5=0$ bằng
	\choice 
		{\True  $\dfrac{5}{3}$}
		 {$\dfrac{7}{3}$}
		{$5$} 
		{$\dfrac{5}{9}$} 
	\loigiai{ 
		Ta có $M(10;0;0)\in(P)$.\\ 
		Vì $\dfrac{1}{1}=\dfrac{2}{2}=\dfrac{2}{2}\ne\dfrac{-10}{-5}$ nên hai mặt phẳng $(P)$ và  $(Q)$ song song.\\
		Khi đó, $\mathrm{d}\Big((P),(Q)\Big)=\mathrm{d}\Big(M,(Q)\Big)=\dfrac{\left|10-5\right|}{\sqrt{9}}=\dfrac{5}{3}$. 
	 }
\end{ex}

% nguồn: Phan1_Lop_12_Hinh_hoc_Chuong_5
\begin{ex}%[2H5H1-5]%[ma:2e92897f]
	Trong không gian $Oxyz$, cho điểm $M(1;0;1)$ và mặt phẳng $(P)\colon 2x+y+2z+5=0$. Khoảng cách từ điểm $M$ đến mặt phẳng  $(P)$ là
	 \choice
		 {$3\sqrt{2}$}
		{\True $3$} 
		{$\dfrac{9\sqrt{2}}{2}$} 
		 {$\sqrt{3}$}
	 \loigiai{
		Ta có $\mathrm{d}\Big(M,(P)\Big)=\dfrac{\left|2\cdot 1+0+2\cdot 1+5\right|}{\sqrt{2^2+1^2+2^2}}=\dfrac{9}{3}=3$. 
	 }
\end{ex}

% nguồn: Phan1_Lop_12_Hinh_hoc_Chuong_5
\begin{ex}%[2H5H1-5]%[ma:b377d612]
	Trong không gian với hệ trục tọa độ $Oxyz$, cho mặt phẳng $(P)\colon 2x-y+2z+1=0$ và mặt phẳng  $(Q)\colon 2x-y+2z+4=0$. Khoảng cách giữa hai mặt phẳng đã cho bằng
	 \choice
		 {\True $1$}
		{$\dfrac{1}{3}$} 
		 {$3$}
		 {$\dfrac{1}{5}$}
	\loigiai{ 
		Dễ thấy $(P) \parallel(Q) \Rightarrow \mathrm{d}((P),(Q))=\mathrm{d}(M,(Q))$ với $M(0; 1; 0) \in(P)$.\\ 
		 Vậy $\mathrm{d}((P),(Q))=\dfrac{|2\cdot 0-1+2\cdot 0+4|}{\sqrt{2^2+(-1)^2+2^2}}=1$.
	 }
\end{ex}

% nguồn: Phan1_Lop_12_Hinh_hoc_Chuong_5
\begin{ex}%[2H5H1-5]%[ma:0bc4d0c0]
	Trong không gian với hệ trục tọa độ $Oxyz$, cho hai mặt phẳng song song $(P)\colon  x+y+z-2=0$; $(Q)\colon x+y+z+4=0$. Khoảng cách giữa hai mặt phẳng $(P)$ và $(Q)$ bằng
	\choice 
		 {$\dfrac{2\sqrt{3}}{3}$}
		{$\sqrt{3}$} 
		{$6$} 
		{\True $2\sqrt{3}$} 
	 \loigiai{
		Ta  có $\mathrm{d}\Big((P),(Q)\Big)=\dfrac{\left|4+2\right|}{\sqrt{1^2+1^2+1^2}}=\dfrac{6}{\sqrt{3}}=2\sqrt{3}$.
	 }
\end{ex}

% nguồn: Phan1_Lop_12_Hinh_hoc_Chuong_5
\begin{ex}%[2H5H1-5]%[ma:9f905572]
	Trong không gian với hệ tọa độ $Oxyz$, mặt phẳng $(\alpha)$ qua hai điểm $C(2;1;6)$, $D(3;5;1)$ và cách đều hai điểm  $A(6;4;0)$, $B(4;5;0)$ có phương trình là
	\choice 
		{$13x + 18y  + 15z + 134 = 0; 5x + 10y + 9z - 7 = 0$}
		 {$13x + 18y + 15z = 0; 5x + 10y + 9z - 4 = 0$}
		{\True $5x + 10y  + 9z - 74 = 0; 13x + 18y + 17z - 146 = 0$}
		{$13x + 18y + 15z - 13 = 0; 5x + 10y + 9z + 4  = 0$}
	\loigiai{ 
		Mặt phẳng $(\alpha)$ cách đều hai điểm $A, B$ khi và chỉ khi $AB \parallel  (\alpha)$ hoặc $(\alpha)$ đi qua trung điểm $I$ của $AB$.\\
		\textbf{Trường  hợp 1:} $(\alpha) \parallel AB$ ta có $\overrightarrow{CD} = (1;4;-5)$, $\overrightarrow{AB} = (-2;1;0)$.\\
		Mặt phẳng $(\alpha)$ qua điểm $C(2;1;6)$ và có vectơ  pháp tuyến là $\overrightarrow{n} = [\overrightarrow{CD}, \overrightarrow{AB}] = (5;10;9)$.\\
		Phương trình mặt phẳng  $(\alpha)$ là $5(x - 2) + 10(y - 1) + 9(z - 6) = 0$ hay $5x + 10y + 9z - 74 = 0$.\\
		\textbf{Trường hợp 2:}  $(\alpha)$ qua trung điểm $I\left(5;\dfrac{9}{2};0\right)$ của $AB$. Có $\overrightarrow{CD} = (1;4;-5)$, $\overrightarrow{CI} =\left(3; \dfrac{7}{2}; -6\right)$.\\
		Mặt phẳng $(\alpha)$ qua điểm $C(2;1;6)$ và  có một vectơ pháp tuyến là
		$$\overrightarrow{n} = \left[\overrightarrow{CD}, \overrightarrow{CI}\right] =  \left(\dfrac{13}{2}; -9; -\dfrac{17}{2}\right).$$
		Chọn một vectơ  pháp tuyến của mặt phẳng $(\alpha)$ là $\overrightarrow{n_{(\alpha)}} = (13;18;17)$.\\
		Phương trình mặt phẳng $(\alpha)$ $$13(x - 2) + 18(y - 1) + 17(z - 6) = 0 \text{ hay } 13x  + 18y + 17z - 146 = 0.$$
	} 
\end{ex}

% nguồn: Phan1_Lop_12_Hinh_hoc_Chuong_5
\begin{ex}%[2H5H1-5]%[ma:01da901c]
	Trong không gian với hệ trục tọa độ $Oxyz$, khoảng cách từ điểm $M(0;3;-1)$  đến mặt phẳng $(\alpha)\colon 2x+y-2z-2=0$ bằng
	\choice 
		{\True $1$} 
		{$\dfrac{4}{3}$} 
		 {$\dfrac{1}{3}$}
		{$3$} 
	 \loigiai{
		Ta  có $\mathrm{d}\Big(M;(\alpha)\Big)=\dfrac{\left|2\cdot 0+3-2\cdot (-1)-2\right|}{\sqrt{2^2+1^2+(-2)^2}}=1$.
	 }
\end{ex}

% nguồn: Phan1_Lop_12_Hinh_hoc_Chuong_5
\begin{ex}%[2H5H1-5]%[ma:a8955909]
	Trong không gian $Oxyz$, khoảng cách từ gốc tọa  độ $O$ đến mặt phẳng $(P)\colon 2x-y+2z+12=0$ bằng
	 \choice
		{$12$} 
		 {$1$}
		 {$\dfrac{4}{3}$}
		 {\True $4$}
	 \loigiai{
		Ta  có $\mathrm{d}\Big(O,(P)\Big)=\dfrac{\left|2\cdot 0-0+2\cdot 0+12\right|}{\sqrt{2^2+(-1)^2+2^2}}=4$.
	} 
\end{ex}

% nguồn: Phan1_Lop_12_Hinh_hoc_Chuong_5
\begin{ex}%[2H5H1-5]%[ma:1f672d1c]
	Trong không gian $(Oxyz)$, cho $A(2;0;0)$, $B(0;4;0)$, $C(0;0;6)$, $D(2;4;6)$. Gọi $(P)$ là mặt phẳng song song với  $(ABC)$, $(P)$ cách đều $D$ và mặt phẳng $(ABC)$. Phương trình của $(P)$ là
	 \choice
		 {\True $6x+3y+2z-24=0$}
		{$6x+3y+2z-12=0$} 
		{$6x+3y+2z=0$} 
		{$6x+3y+2z-36=0$} 
	\loigiai{ 
		 $(ABC)\colon \dfrac{x}{2}+\dfrac{y}{4}+\dfrac{z}{6}=1 \Leftrightarrow 6x+3y+2z-12=0$.\\
		$(P)\parallel (ABC) \Rightarrow(P)\colon  6x+3y+2z+m=0$, $(m \neq-12)$.\\
		$(P)$ cách đều $D$  và mặt phẳng $(ABC)$ nên
		\begin{eqnarray*}
			\mathrm{d}\left(D,(P)\right)=\mathrm{d}\left(D,(ABC)\right)&\Leftrightarrow&\dfrac{\left|6\cdot2+3\cdot4+2\cdot6+m\right|}{\sqrt{6^2+3^2+2^2}}=\dfrac{\left|6\cdot2+3\cdot0+2\cdot0+m\right|}{\sqrt{6^2+3^2+2^2}}\\ 
			& \Leftrightarrow&\left|36+m\right|=\left|12+m\right|\\  &\Leftrightarrow&\hoac{&36+m=12+m\\ &36+m=-12-m}\\
			&\Leftrightarrow&  m=-24.
		\end{eqnarray*}
		Vậy phương trình  của $(P)$ là $6x+3y+2z-24=0$.
	 }
\end{ex}

% nguồn: Phan1_Lop_12_Hinh_hoc_Chuong_5
\begin{ex}%[2H5H1-5]%[ma:6ecac7f7]
	Cho mặt phẳng $(P)\colon 3 x+4 y+2  z+4=0$ và điểm $A(1;-2; 3)$. Khoảng cách từ $A$ đến $(P)$ bằng
	 \choice
		{\True $\dfrac{5}{\sqrt{29}}$} 
		 {$\dfrac{5}{29}$}
		 {$\dfrac{\sqrt{5}}{3}$}
		{$\dfrac{5}{9}$} 
	\loigiai{ 
		Khoảng  cách từ $A$ đến $(P)$ là
		$$\mathrm{d}(A,(P))=\dfrac{\left|  3\cdot 1+ 4\cdot (-2)+2\cdot 3+4\right| }{\sqrt{3^2+4^2+2^2}}=\dfrac{5}{\sqrt{29}}.$$
	} 
\end{ex}

% nguồn: Phan1_Lop_12_Hinh_hoc_Chuong_5
\begin{ex}%[2H5H1-5]%[ma:c74fcacd]
	Trong  không gian $Oxyz$, cho mặt phẳng $(P)\colon x+y+4z-3=0$ và điểm $A(1;1;3)$. Mặt phẳng $(Q)\parallel(P)$ và cắt các tia $Ox$, $Oy$ lần lượt tại các điểm $B$ và $C$ sao cho tam giác $ABC$ có diện tích bằng $2\sqrt{22}$. Khoảng cách từ điểm $M(2;2;1)$ đến $(Q)$ bằng
	\choice 
		{\True  $2\sqrt{2}$}
		 {$\dfrac{8\sqrt{6}}{3}$}
		{$\dfrac{\sqrt{6}}{3}$} 
		 {$\dfrac{2\sqrt{2}}{3}$}
	\loigiai{ 
		Mặt phẳng $(Q)\parallel(P)\Rightarrow  (Q)$ có dạng $x+y+4z+d=0$, $(d\ne-3)$.\\
		$(Q)\cap Ox=B(-d;0;0)$, $(Q)\cap  Oy=C(0;-d;0)$.\\
		Do $B$, $C$ lần lượt thuộc các tia $Ox$,  $Oy\Rightarrow d<0$.\\
		Ta có $\overrightarrow{AB}=(-d-1;-1;-3)$, $\overrightarrow{AC}=(-1;-d-1;-3)\Rightarrow  \left[\overrightarrow{AB},\overrightarrow{AC}\right]=(-3d;-3d;d^2+2d)$.\\
		 $S_{\triangle ABC}=2\sqrt{22}\Leftrightarrow\dfrac{1}{2}\left|\left[\overrightarrow{AB},\overrightarrow{AC}\right]\right|=2\sqrt{22}\Leftrightarrow 9d^2+9d^2+(d^2+2d)^2=352$.\\
		Suy ra $d^4+4d^3+22d^2-352=0\Leftrightarrow  (d+4)(d^3+22d-88)=0\Leftrightarrow \hoac{&d+4=0\\&d^3+22d-88=0.}$\\
		\begin{itemize}
			\item Với $d+4=0\Leftrightarrow d=-4$ (thỏa  mãn).
			\item Với  $d^3+22d-88=0$\\
			Do $d<0$ nên ta có $d^3+22d-88<0$,  $\forall d\in (-\infty;0)$.\\
			Suy ra phương trình  $d^3+22d-88=0$ không có nghiệm âm.
		\end{itemize}
		 Suy ra chỉ có $d=-4$ thỏa mãn.\\
		Khi đó ta có phương trình mặt  phẳng $(Q)\colon x+y+4z+4=0$.\\
		 Khoảng cách từ điểm $M(2;2;1)$ đến $(Q)$ bằng $\dfrac{\left|2+2+4\cdot 1+4\right|}{\sqrt{1^2+1^2+4^2}}=2\sqrt{2}$.
	} 
\end{ex}

% nguồn: Phan1_Lop_12_Hinh_hoc_Chuong_5
\begin{ex}%[2H5H1-5]%[ma:0629ee25]
	Trong  không gian với hệ tọa độ $Oxyz$, mặt phẳng $(Q)$ song song với mặt phẳng $(P)\colon x + 2y - 2z + 1 = 0$ và $(Q)$ cách điểm $M(1;-2;1)$ một khoảng bằng $3$ có phương trình là
	 \choice
		{$x + 2y - 2z + 10  = 0$; $x + 2y - 2z - 1 = 0$}
		{$x + 2y - 2z  = 0$; $x + 2y - 2z - 14 = 0$}
		{\True $x +  2y - 2z + 14 = 0$; $x + 2y - 2z - 4 = 0$}
		{$x + 2y - 2z + 4 = 0$; $x + 2y - 2z -  14 = 0$}
	\loigiai{ 
		Mặt phẳng $(Q)$ song song với mặt phẳng $(P)\colon x + 2y - 2z  + 1 = 0$ nên mặt phẳng $(Q)$ có phương trình $x + 2y - 2z + D = 0$\, $(D \neq 1)$.\\
		Mặt phẳng $(Q)$ cách điểm $M(1;-2;1)$ một khoảng bằng $3$  nên
		\begin{eqnarray*}
			&&\mathrm{d}(M;(Q)) = 3\\ 
			&&  \Leftrightarrow \dfrac{|1 + 2(-2) - 2(1) + D|}{\sqrt{1^2 + 2^2 + (-2)^2}} = 3\\
			&& \Leftrightarrow \dfrac{|1 - 4 - 2 + D|}{3} =  3\\
			&& \Leftrightarrow |D - 5| = 9\\ 
			 && \Leftrightarrow \hoac{&D-5=9\\&D-5=-9}\\
			&& \Leftrightarrow \hoac{&D=14\\&D=-4} (\text{thỏa mãn}). 
		\end{eqnarray*}
		Vậy có hai mặt  phẳng $(Q)$ là $x + 2y - 2z + 14 = 0$ hoặc $x + 2y - 2z - 4 = 0$.
	 }
\end{ex}

% nguồn: Phan1_Lop_12_Hinh_hoc_Chuong_5
\begin{ex}%[2H5H1-5]%[ma:2659b0cb]
	 Trong không gian $Oxyz$, cho điểm $A(0;1;2)$ và đường thẳng $d\colon \dfrac{x-2}{2}=\dfrac{y-1}{2}=\dfrac{z-1}{-3}$. Gọi $(P)$ là mặt phẳng đi qua $A$ và chứa $d$. Khoảng cách từ điểm $M(5;-1;3)$ đến $(P)$ bằng
	 \choice
		 {$5$}
		{$\dfrac{1}{3}$} 
		 {\True $1$}
		{$\dfrac{11}{3}$} 
	 \loigiai{
		Chọn $B(2;1;1)\in  d$ ta có $\overrightarrow{AB}=(2;0;-1)$ và có $\left[\overrightarrow{AB},\overrightarrow{u}_d\right]=(2;4;4)=2(1;2;2)$.\\
		Mặt phẳng $(P)$  đi qua $A$ và chứa $d$ suy ra $\overrightarrow{n}_P=(1;2;2)\Rightarrow (P)\colon x+2y+2z-6=0$.\\
		 Vậy $\mathrm{d}\Big(M,(P)\Big)=\dfrac{\left|x_M+2y_M+2z_M-6\right|}{\sqrt{1^2+2^2+2^2}}=1$.
	 }
\end{ex}

% nguồn: Phan1_Lop_12_Hinh_hoc_Chuong_5
\begin{ex}%[2H5H1-5]%[ma:8c8ad44d]
	Trong không gian $Oxyz$, cho điểm $M(2;-1;3)$ và mặt phẳng $(P)\colon 2x-2y+z+1=0$. Khoảng cách điểm $M$ đến mặt phẳng $(P)$ bằng 
	\choice 
		{$2$} 
		{$\dfrac{5}{3}$} 
		{$3$} 
		 {\True $\dfrac{10}{3}$}
	 \loigiai{
		Ta  có $\mathrm{d}\Big(M,(P)\Big)=\dfrac{\left|2\cdot 2-2\cdot (-1)+3+1\right|}{\sqrt{2^2+(-2)^2+1^2}}=\dfrac{10}{3}$.
	} 
\end{ex}

% nguồn: Phan1_Lop_12_Hinh_hoc_Chuong_5
\begin{ex}%[2H5H1-5]%[ma:ffb2561b]
	Trong không gian với hệ trục $Oxyz$, cho mặt phẳng $(P)\colon x - 2y + z - 1 = 0$ và mặt phẳng $(Q)\colon -x + 2y - z + 7 = 0$. Khi đó khoảng cách giữa hai mặt phẳng  $(P)$ và $(Q)$ bằng
	 \choice
		 {$1$}
		{$\dfrac{4\sqrt{6}}{3}$} 
		 {$\dfrac{3}{2}$}
		{\True $\sqrt{6}$} 
	 \loigiai{
		Ta nhận thấy $(Q)$  có thể viết lại thành
		 \[
		Q\colon -(x - 2y + z) + 7 = 0 \Leftrightarrow x -  2y + z - 7 = 0.
		 \]
		So với $(P)\colon x - 2y + z - 1 = 0$, ta thấy hai mặt phẳng song  song.\\
		Khoảng cách giữa hai mặt phẳng  song song
		\[ 
		\mathrm{d} = \dfrac{|d_1 - d_2|}{\sqrt{a^2  + b^2 + c^2}} = \dfrac{| -1 - (-7) |}{\sqrt{1^2 + (-2)^2 + 1^2}} = \dfrac{6}{\sqrt{6}} = \sqrt{6}.
		 \]
	} 
\end{ex}

% nguồn: Phan1_Lop_12_Hinh_hoc_Chuong_5
\begin{ex}%[2H5H1-5]%[ma:cffcaf74]
	Trong không gian hệ trục tọa độ $Oxyz$, khoảng cách từ điểm  $A(3;1;-2)$ đến mặt phẳng $z=0$ bằng
	 \choice
		{$\sqrt{5}$} 
		 {$\sqrt{14}$}
		 {\True $2$}
		{$3$} 
	\loigiai{ 
		Khoảng cách từ $A(3;1;-2)$ đến  mặt phẳng $z=0$ bằng $2$.
	} 
\end{ex}

% nguồn: Phan1_Lop_12_Hinh_hoc_Chuong_5
\begin{ex}%[2H5H1-5]%[ma:698ed9c1]
	Trong không gian với hệ trục tọa độ $Oxyz$, mặt phẳng $(\alpha)$ cắt các trục $Ox$, $Oy$, $Oz$ lần lượt  tại $3$ điểm $A(2;0;0)$, $B(0;3;0)$, $C(0;0;-4)$. Khoảng cách từ $O$ đến $(\alpha)$ bằng
	\choice 
		{$\dfrac{\sqrt{61}}{12}$} 
		 {$4$}
		{\True  $\dfrac{12\sqrt{61}}{61}$}
		 {$3$}
	 \loigiai{
		Phương trình mặt phẳng $(ABC)$ có dạng  $\dfrac{x}{2}+\dfrac{y}{3}+\dfrac{z}{-4}=1\Leftrightarrow 6x+4y-3z-12=0$.\\
		Khoảng cách từ $O$ đến $(\alpha)$ bằng $\mathrm{d}\Big(O,(\alpha)\Big)=\dfrac{\left|6\cdot 0+4\cdot 0-3\cdot 0-12\right|}{\sqrt{6^2+4^2+(-3)^2}}=\dfrac{12\sqrt{61}}{61}$. 
	 }
\end{ex}

% nguồn: Phan1_Lop_12_Hinh_hoc_Chuong_5
\begin{ex}%[2H5H1-5]%[ma:57180ba5]
	Trong không gian với hệ tọa độ $Oxyz$, cho điểm $B(1;-1;0)$, $C(2;-1;-1)$, $M(-2;1;3)$. Mặt phẳng $(P)$ vuông góc với $BC$ và cách $M$  một khoảng bằng $\sqrt{2}$ có phương trình là
	 \choice
		{$x - z -  7 = 0$; $x - z - 3 = 0$}
		{\True $x -  z + 7 = 0$; $x - z + 3 = 0$}
		{$x - z + 7 =  0$; $x - z = 0$}
		{$x - z - 7 = 0$; $x - z  + 3 = 0$}
	\loigiai{ 
		Mặt phẳng $(P)$ vuông góc với $BC$ nên  $\overrightarrow{n_{(P)}} = \overrightarrow{BC} = (1;0;-1) \Rightarrow (P)\colon x - z + D = 0$.\\
		\begin{eqnarray*}
			&&\mathrm{d}(M;(Q)) =  \sqrt{2}\\
			&& \Leftrightarrow \dfrac{|-2 - 3 +  D|}{\sqrt{1^2 + (-1)^2}} = \sqrt{2}\\
			&& \Leftrightarrow  |D - 5| = 2\\
			&&  \Leftrightarrow \hoac{&D-5=2\\&D-5=-2}\\
			 && \Leftrightarrow \hoac{&D=7\\&D=3} (\text{thỏa mãn}).
		\end{eqnarray*}
		Vậy có hai mặt phẳng $(P)$ là $x - z + 7 =  0$ hoặc $x - z + 3 = 0$.
	 }
\end{ex}

% nguồn: Phan1_Lop_12_Hinh_hoc_Chuong_5
\begin{ex}%[2H5H1-5]%[ma:35a5f5d5]
	 Trong không gian $Oxyz$, khoảng cách từ điểm $A(1;0;0)$ tới mặt phẳng $(P)\colon 2x+2y-z+1=0$ bằng
	\choice 
		{$3$} 
		 {$\sqrt{3}$}
		 {$9$}
		{\True  $1$}
	\loigiai{Khoảng cách từ điểm $A(1;0;0)$ tới mặt phẳng  $(P)\colon 2x+2y-z+1=0$ là $$\mathrm{d}(A,(P))=\dfrac{|2\cdot1+2\cdot0-0+1|}{\sqrt{2^2+2^2+(-1)^2}}=\dfrac{|3|}{\sqrt{9}}=1.$$}
\end{ex}

% nguồn: Phan1_Lop_12_Hinh_hoc_Chuong_5
\begin{ex}%[2H5H1-5]%[ma:499a2f0c]
	Trong không gian với hệ trục tọa độ $Oxyz$, cho hai mặt phẳng  song song $(P)\colon 2x+y-2z-1=0$, $(Q)\colon 6x+3y-6z+15=0$. Khoảng cách giữa hai mặt phẳng $(P)$, $(Q)$ bằng
	 \choice
		 {\True $2$}
		{$\dfrac{4}{3}$} 
		 {$\dfrac{16}{9}$}
		{$\dfrac{16}{3}$} 
	\loigiai{ 
		Ta có $(Q)\colon 6x+3y-6z+15=0\Leftrightarrow 2x+y-2z+5=0$.\\ 
		Lấy  $M(0;1;0)\in(P)$.\\
		Vậy $\mathrm{d}\Big((P),(Q)\Big)=\mathrm{d}\Big(M,(Q)\Big)=\dfrac{\left|1+5\right|}{\sqrt{1+4+4}}=2$. 
	} 
\end{ex}

% nguồn: Phan1_Lop_12_Hinh_hoc_Chuong_5
\begin{ex}%[2H5H1-5]%[ma:02a9a33a]
	Trong không gian với hệ tọa độ $Oxyz$, cho  hai mặt phẳng song song $(P)\colon x+y+z-2=0$; $(Q)\colon x+y+z+4=0$. Khoảng cách giữa hai mặt phẳng $(P)$ và $(Q)$ bằng
	 \choice
		{$\dfrac{2\sqrt{3}}{3}$} 
		 {$\sqrt{3}$}
		 {$6$}
		{\True $2\sqrt{3}$} 
	 \loigiai{
		 Lấy điểm $A(0;0;-4)\in(Q)$.\\
		Vì $(P)\parallel (Q)$  nên
		\[\mathrm{d}((P),(Q))=\mathrm{d}(A,(P))=\dfrac{|0+0-4-2|}{\sqrt{1^2+1^2+1^2}}=\dfrac{6}{\sqrt{3}}=2\sqrt{3}. \] 
	} 
\end{ex}

% nguồn: Phan1_Lop_12_Hinh_hoc_Chuong_5
\begin{ex}%[2H5V1-5]%[ma:bf47726e]
	Trong không gian $Oxyz$, cho ba điểm $A(1;0;0), B(0;-2;0), C(0;0; 1)$. Khoảng cách từ gốc  toạ độ đến mặt phẳng $(ABC)$ bằng
	 \choice
		 {$\dfrac{1}{3}$}
		 {$-\dfrac{2}{3}$}
		{\True $\dfrac{2}{3}$} 
		{$\dfrac{2}{\sqrt{7}}$} 
	\loigiai{ 
		Ta có phương trình mặt phẳng qua  đoạn chắn $(ABC)$
		$$(ABC) \colon  \dfrac{x}{1}+\dfrac{y}{-2}+\dfrac{z}{1}=1 \Leftrightarrow -2x+y-2z+2=0$$
		Vậy khoảng cách từ  gốc tọa~độ $O$ đến mặt phẳng $(ABC)$~là
		 $$\mathrm{d}\left(O;(ABC)\right)=\dfrac{\left|-2\cdot0+0-2\cdot0+2\right|}{\sqrt{(-2)^2+1^2+(-2)^2}}=\dfrac{2}{3}$$
	 }
\end{ex}


%%%%%%%%%%%%%%%% ĐÚNG/SAI %%%%%%%%%%%%%%%%

% nguồn: Phan1_Lop_12_Hinh_hoc_Chuong_5
\begin{ex}%[2H5H1-5]%[ma:77479aed]
	Trong không gian $Oxyz$, cho mặt phẳng $(P)$  có phương trình $10x + 7y - 26z + 2 = 0$.
	\choiceTF 
		{\True $\overrightarrow{n} = \left( 10; 7; -26 \right)$ là một vectơ pháp tuyến của mặt phẳng $(P)$} 
		{\True Khoảng cách từ  điểm $M\left( 0;1;-1 \right)$ đến mặt phẳng $(P)$ bằng $\dfrac{7\sqrt{33}}{33}$}
		{Mặt phẳng $(P)$ vuông  góc với mặt phẳng $(Q)\colon 10x + 7y - 26z + 2 = 0$}
		{Mặt phẳng $(\alpha)$ đi qua hai điểm $A\left( -2;-2;-2 \right)$, $B\left( 0;-2;-1 \right)$ và vuông góc với mặt phẳng $(P)$  có phương trình dạng $ax + by + cz - 138 = 0$. Khi đó $a + b + c = -69$}
	 \loigiai{
		\begin{itemchoice}
			\itemch $\overrightarrow{n} =  \left( 10; 7; -26 \right)$ là một vectơ pháp tuyến của mặt phẳng $(P)$.
			\itemch $\mathrm{d}(M;(P))=\dfrac{|7+26+2|}{\sqrt{10^2+7^2+26^2}}=\dfrac{7\sqrt{33}}{33}$. 
			\itemch Dễ thấy mặt phẳng $(P)$ trùng mặt  phẳng $(Q)$.
			\itemch Ta có $\overrightarrow{AB}=(2;0;1)$.\\ 
			Vì mặt phẳng $(\alpha)$ đi qua $A$, $B$ đồng thời vuông góc với $(P)$ nên mặt phẳng $(\alpha)$ có véc-tơ pháp  tuyến là $\overrightarrow{n_\alpha}=\left[\overrightarrow{AB};\overrightarrow{n}\right]=(33;-62;-14)$\\
			Suy ra phương trình mặt phẳng $(\alpha)  \colon 33x-66y-14z-138=0$.\\
			Vậy $a + b +  c = -43$.
		\end{itemchoice}
	 }
\end{ex}

% nguồn: Phan1_Lop_12_Hinh_hoc_Chuong_5
\begin{ex}%[2H5H1-5]%[2H5H3-3]%[ma:95899593]
	Trong không gian với hệ tọa  độ $Oxyz$, cho mặt phẳng $(P)\colon x-2y+2z-2=0$ và điểm $I(-1; 2;-1)$. Biết mặt cầu $(S)$ có tâm $I$ và cắt mặt phẳng $(P)$ theo giao tuyến là đường tròn $(C.$ có diện tích là $25\pi$. Xét tính đúng sai của các khẳng định sau
	 \choiceTF
		{\True  Bán kính đường tròn $(C)$ là $r=5$}
		{\True Khoảng cách từ $I$  đến mặt phẳng $(P)$ là $3$}
		{Tâm đường tròn $(C)$ có tọa độ là $H(1;  3; 1)$}
		{Phương trình  mặt cầu $(S)$ là $(x+1)^2+(y-2)^2+(z+1)^2=16$}
	\loigiai{ 
		Gọi $M$ là điểm nằm trên đường tròn giao tuyến  của $(S)$ và $(P)$.\\
		Ta có $IM=R$. Áp dụng công thức tính bán kính mặt cầu trong trường hợp mặt cầu $(S)$ giao với mặt phẳng $(P)$ theo giao  tuyến là đường tròn có bán kính $r$ là
		\begin{itemchoice}
			\itemch  Ta có đường tròn ($C$) có diện tích là $25\pi$. Suy ra $r=5=MH$.
			 \itemch Ta có $\mathrm{d}(I;(P))=\dfrac{|-1-2.2+2.(-1)-2|}{\sqrt{1+(-2)^2+2^2}}=3=IH$
			\itemch Phương trình $IH$: $\heva{&x=-1+t \\&y=2-2t \\ &z=-1+2t.}$\\ 
			Vậy tọa độ là $H$ là nghiệm của hệ $\left\{\begin{array}{l}x=-1+t \\ y=2-2t \\ z=-1+2t \\ x-2y+2z-2=0
		\end{array} \Leftrightarrow\left\{\begin{array}{l}x=0\\ y=0\\ z=1\\ t=1
	\end{array} \Rightarrow H(0; 0; 1)\right.\right.$.
	\itemch Suy ra  $R^2=MI^2=IH^2+MH^2=3^2+5^2=34$.\\
	Vậy phương trình mặt  cầu $(S)$ là $(x+1)^2+(y-2)^2+(z+1)^2=34$.
	\end{itemchoice}
	} 
	\end{ex}

% nguồn: Phan1_Lop_12_Hinh_hoc_Chuong_5
\begin{ex}%[2H5H1-5]%[ma:ce09b842]
	Trong không gian $Oxyz$, cho  các điểm $A(2;-3;1)$, $B(1;0;-1)$, $C(3;1;2)$, $D(1;0;-1)$.
	 \choiceTF
		{Mặt phẳng đi qua điểm $A$ và song song với  mặt phẳng $(Oxz)$ có phương trình là $y - 3 = 0$}
		{\True Mặt phẳng đi  qua điểm $A$ và song song với mặt phẳng trung trực của $BC$ có phương trình là $2x + y + 3z - 4 = 0$}
		{\True Gọi $(\alpha)$ là mặt phẳng đi qua trọng tâm của $\triangle ABC$ và song song với mp $(\beta)\colon z - 2 = 0$. Khoảng cách  từ điểm $A$ đến $(\alpha)$ bằng $ \dfrac{1}{3}$}
		{Mặt phẳng đi qua điểm $D$ và  song song với mặt phẳng $(ABC)$ có phương trình là $11x - y - 7z - 18 = 0$}
	\loigiai{ 
		\begin{itemchoice}
			\itemch  \textbf{Sai}. Gọi $(P)$ là mp cần tìm. Do $(P) \parallel (Oxz) \Rightarrow \overrightarrow{n}_P = \overrightarrow{j} = (0;1;0)$.\\
			Và $(P)$ đi qua điểm $A(2;-3;1) \Rightarrow (P) \colon 0\cdot (x-2) +  1\cdot (y+3) + 0 \cdot (z-1) = 0 \Rightarrow (P)\colon y+3=0$.
			\itemch \textbf{Đúng}. Trung điểm $BC$ là $I\left(2; \dfrac{1}{2}; \dfrac{1}{2}\right)$, $\overrightarrow{BC} = (2;1;3)$.\\ 
			Mặt  phẳng trung trực của $BC$ là $(Q)\colon 2 \cdot (x-2) + 1\cdot \left(y- \dfrac{1}{2}\right) + 3 \cdot \left(z- \dfrac{1}{2}\right) = 0$\\
			$\Rightarrow (Q)\colon 2x +  y + 3z - 6 = 0$.\\
			Gọi $(R)$ là mặt phẳng cần tìm. Do  $(R)\parallel(Q) \Rightarrow (R)\colon 2 \cdot (x-2) + 1 \cdot (y+3) + 3\cdot (z-1) = 0$\\
			 $\Rightarrow (R)\colon 2x + y + 3z - 4 = 0$.
			 \itemch \textbf{Đúng}. Trọng tâm $\triangle ABC$ là điểm $G\left(2;- \dfrac{2}{3}; \dfrac{2}{3}\right)$.\\
			Do $(\alpha)\parallel  (\beta) \Rightarrow \overrightarrow{n}_{\alpha} = \overrightarrow{n}_{\beta} = (0;0;1)$ và $(\alpha)$ đi qua điểm $G$\\
			$\Rightarrow (\alpha)\colon 0\cdot(x-2) + 0 \cdot \left(y+ \dfrac{2}{3}\right) + 1 \cdot \left(z- \dfrac{2}{3}\right) =  0 \Rightarrow (\alpha)\colon 3z-2=0$ và\\
			$d(A;(\alpha)) = \dfrac{|3\cdot  1-2|}{\sqrt{0^2+0^2+3^2}} = \dfrac{1}{3}$.
			\itemch \textbf{Sai}. Ta có $\overrightarrow{AB}=(-1;3;-2), \overrightarrow{AC}=(1;4;1)$. Suy ra vectơ pháp tuyến là $\overrightarrow{n}_{ABC} =  [\overrightarrow{AB}, \overrightarrow{AC}] = (11;-1;-7)$.\\
			Gọi $(Q)$ là mặt phẳng cần tìm. Do $(Q)\parallel (ABC)$ nên vectơ pháp tuyến $\overrightarrow{n}_{Q} = \overrightarrow{n}_{ABC} = (11;-1;-7)$ và $(Q)$ đi qua điểm $D(1;0;-1) \Rightarrow (Q)\colon 11\cdot (x-1)-1 \cdot (y-0)-7  \cdot (z+1)=0 \Rightarrow (Q)\colon 11x-y-7z-18=0$.\\
			Nhận  thấy tọa độ điểm $A$ thỏa mãn phương trình $(Q)$ (vì $11 \cdot 2 - 1 \cdot (-3) - 7 \cdot 1 - 18 = 0$) nên mặt phẳng $(Q)$ chính là mặt phẳng $(ABC)$.
		\end{itemchoice}
	 }
\end{ex}

% nguồn: Phan1_Lop_12_Hinh_hoc_Chuong_5
\begin{ex}%[2H5H1-5]%[ma:542825b2]
	Trong không gian với hệ toạ độ $Oxyz$, cho mặt phẳng $\left(  P \right)$ đi qua điểm $M\left( -1;2;0 \right)$ và có một vectơ pháp tuyến $\overrightarrow{n}=\left( 4;0;-5 \right)$.
	\choiceTF 
		{\True Mặt phẳng $\left( P \right)$ có  phương trình là $4x-5z+4=0$}
		 {Mặt phẳng $\left( P \right)$ đi qua điểm $A\left( 2;-1;5 \right)$}
		{Mặt phẳng $\left( Q \right)\colon  4x-5z+4=0$ song song với $\left( P \right)$}
		{\True Cho điểm $B\left( -2;1;3 \right)$. Khi đó $\mathrm{d}\left( A,\left( P \right)  \right)=\dfrac{19}{\sqrt{41}}$}
	 \loigiai{
		\begin{itemchoice}
			 \itemch \textbf{Đúng}. Phương trình mặt phẳng $\left( P \right)\colon 4\left( x+1 \right)+0\left( y-2 \right)-5\left( z-0 \right)=0\Leftrightarrow 4x-5z+4=0$.
			\itemch \textbf{Sai}. Thay điểm $M$ vào phương trình mặt phẳng $\left( P \right)\colon 4x-5z+4=0$ ta  thấy không thỏa mãn.
			\itemch \textbf{Sai}. Do $\left( Q \right)\equiv \left( P \right)$. 
			\itemch \textbf{Đúng}. Do $\mathrm{d}\left( B,\left( P \right) \right)=\dfrac{\left| 4 \cdot (-2)-5\cdot 3+4  \right|}{\sqrt{{{4}^{2}}+{{5}^{2}}}}=\dfrac{19}{\sqrt{41}}$.
		\end{itemchoice}
	} 
\end{ex}

% nguồn: Phan1_Lop_12_Hinh_hoc_Chuong_5
\begin{ex}%[2H5H1-5]%[ma:b7379909]
	Trong  không gian $Oxyz$, cho mặt phẳng $(P)$ đi qua ba điểm $A(4;0;-3)$, $B(3;2;3)$ và $C(4;-1;-5)$.
	 \choiceTF
		{\True  Mặt phẳng $(P)$ có một cặp véc-tơ chỉ phương là $\overrightarrow{AB}=(-1; 2; 6), \overrightarrow{AC}=(0;-1;-2)$}
		{\True Mặt phẳng  $(P)$ có một véc-tơ pháp tuyến là $\overrightarrow{n}=(2;-2;1)$}
		{\True Mặt phẳng $(P)$ có phương trình  là $2x-2y+z-5=0$}
		{Khoảng cách từ  điểm $M(0;0;11)$ đến mặt phẳng $P$ bằng $3$}
	\loigiai{ 
		\begin{itemchoice}
			\itemch \textbf{Đúng}. Ta có $\overrightarrow{AB}=(-1;2;6)$ và $\overrightarrow{AC}=(0;-1;-2)$.\\ 
			Suy ra mặt phẳng $(P)$ có một cặp véc-tơ chỉ phương là  $\overrightarrow{AB}=(-1; 2; 6), \overrightarrow{AC}=(0;-1;-2)$.
			\itemch \textbf{Đúng}. Ta có $\left[\overrightarrow{AB},\overrightarrow{AC}\right]=(2;-2;1)$.\\ 
			Do đó mặt  phẳng $(P)$ có một véc-tơ pháp tuyến là $\overrightarrow{n}=(2;-2; 1)$.
			\itemch \textbf{Đúng}. Mặt phẳng $(P)$ đi qua điểm $A(4;0;-3)$ có một véc-tơ pháp tuyến là $\overrightarrow{n}=(2;-2;1)$ nên  có phương trình \[2(x-4)-2(y-0)+1(z+3)=0\Leftrightarrow2x-2y+z-5=0.\]
			 \itemch \textbf{Sai}. Khoảng cách từ điểm $M(0;0;11)$ đến mặt phẳng $(P)\colon 2x-2y+z-5=0$ là
			 \[\mathrm{d}\left(M,(P)\right)=\dfrac{|2\cdot0-2\cdot0+11-5|}{\sqrt{2^{2}+(-2)^{2}+1^{2}}}=2.\]
		\end{itemchoice}
	 }
\end{ex}

% nguồn: Phan1_Lop_12_Hinh_hoc_Chuong_5
\begin{ex}%[2H5V1-5]%[ma:b022d135]
	 Trong không gian $Oxyz$, cho hai mặt phẳng $(P)\colon 2x+2y-z+3=0$, $(Q)\colon x+2y-2z+4=0$ và điểm $A(1;2;3)$.
	 \choiceTF
		{\True Điểm $A$ cách mặt phẳng  $(P)$ một khoảng bằng $5$}
		{\True  Mặt phẳng $(Q)$ cắt mặt phẳng $(P)$}
		{Mặt phẳng $(R)\colon 2x+2y-z=0$ cách  mặt phẳng $(P)$ một khoảng bằng $3$}
		{Với mọi  giá trị $m$ thì hai mặt phẳng $(P)$ và $(T)\colon x+y+mz+1=0$ cắt nhau}
	\loigiai{ 
		\begin{itemchoice}
			\itemch Ta có $\mathrm{d}\Big(A,(P)\Big)=\dfrac{\left|2+4-3+3\right|}{\sqrt{2^2+2^2+(-1)^2}}=2$. 
			\itemch Ta có véc-tơ pháp tuyến của hai mặt  phẳng $(P)$ và $(Q)$ lần lượt là $\overrightarrow{n}_P=(2;2;-1)$, $\overrightarrow{n}_Q=(1;2;-2)$.\\
			Hơn nữa $\dfrac{2}{1}\ne\dfrac{2}{2}\ne\dfrac{-1}{-2}$ nên hai véctơ $\overrightarrow{n}_P$ và $\overrightarrow{n}_Q$  không cùng phương.\\
			Do đó hai mặt phẳng $(P)$  và $(Q)$ cắt nhau.
			\itemch Ta có hai mặt phẳng $(P)$ và $(R)$ song song  nhau.\\
			Lấy điểm  $O(0;0;0)$ thuộc mặt phẳng $(R)$ thì $\mathrm{d}\Big((R),(P)\Big)=\mathrm{d}\Big(O,(P)\Big)=\dfrac{\left|3\right|}{\sqrt{2^2+2^2+(-1)^2}}=1$.
			\itemch Mặt phẳng $(P)$ có véc-tơ pháp tuyến là $\overrightarrow{n}_P=(2;2;-1)$.\\ 
			Mặt phẳng $(T)$ có  véc-tơ pháp tuyến là $\overrightarrow{n}_T=(1;1;m)$.\\
			Vì hai mặt phẳng $(P)$ và $(T)$ cắt nhau nên hai véc-tơ pháp tuyến $\overrightarrow{n}_P$  và $\overrightarrow{n}_T$ không cùng phương.\\
			Suy ra $m\ne  -\dfrac{1}{2}$.\\
			Vậy với $m\ne -\dfrac{1}{2}$  thì hai mặt phẳng $(P)$ và $(T)$ cắt nhau.
		\end{itemchoice}
	} 
\end{ex}

% nguồn: Phan1_Lop_12_Hinh_hoc_Chuong_5
\begin{ex}%[2H5V1-5]%[ma:0c0caf10]
	Trong không gian $Oxyz$, cho $ A(-1;0;2)$ và mặt  phẳng $(P)\colon 2x+y-2z=0$.
	 \choiceTF
		{Mặt phẳng $(P)$  có véc-tơ pháp tuyến $\overrightarrow{n}=(2;1;-2)$ và đi qua $A$}
		{\True Phương trình mặt phẳng qua  $A$ và song song $(P)$ là $2x+y-2z+6=0$}
		{\True  Mặt phẳng qua gốc toạ độ $O$, điểm $A$ và vuông góc $(P)$ có phương trình là $2x-2y+z=0$}
		{Có tất cả hai mặt  phẳng song song với $(P)$ và có khoảng cách đến $A$ bằng $2$}
	 \loigiai{
		\begin{itemchoice}
			\itemch Thay toạ độ điểm $A(-1;0;2)$ vào phương trình mặt phẳng $(P)$ ta  được $-2-4=0$ (vô lý).
			\itemch Gọi $(\alpha)$ là mặt phẳng qua $A$ và song song $(P)$, suy  ra $(\alpha)$ có véc-tơ pháp tuyến $\overrightarrow{n}_{\alpha}=(2;1;-2)$.\\
			Phương trình $(\alpha)$ là  $2(x+1)+1(y-0)-2(z-2)=0\Leftrightarrow 2x+y-2z+6=0$.
			\itemch  Gọi $(\beta)$ là mặt phẳng qua $O$, $A$ và vuông góc $(P)$.\\
			Ta có $\overrightarrow{n}_P=(2;1;-2)$, $\overrightarrow{OA}=(-1;0;2)$.\\ 
			Khi đó $(\beta)$ có véc-tơ pháp tuyến  là $\overrightarrow{n}_{\beta}=\left[\overrightarrow{n}_P,\overrightarrow{OA}\right]=(2;-2;1)$.\\
			 Phương trình của $(\beta)$ là $2(x-0)-2(y-0)+1(z-0)=0\Leftrightarrow 2x-2y+z=0$.
			\itemch Phương trình mặt phẳng $(Q)$ có dạng $2x+y-2z+d=0$, với $d\ne  0$.\\
			Ta có 
			 \allowdisplaybreaks
			\begin{eqnarray*}
				 \mathrm{d}\Big(A,(Q)\Big)=2&\Leftrightarrow&\dfrac{\left|2(-1)+0-2\cdot 2+d\right|}{\sqrt{2^2+1^2+(-2)^2}}=2\\
				&\Leftrightarrow&\left|d-6\right|=6\\ 
				 &\Leftrightarrow&\hoac{&d-6=6\\&d-6=-6}\\
				&\Leftrightarrow&\hoac{&d=12&&\text{ (nhận)}\\&d=0.&&\text{ (loại)}} 
			\end{eqnarray*}
		\end{itemchoice}
	} 
\end{ex}

% nguồn: Phan1_Lop_12_Hinh_hoc_Chuong_5
\begin{ex}%[2H5V1-5]%[ma:f40425ca]
	Trong không gian $Oxyz$, cho $A(2;0;0)$, $B(0;2;0)$,  $C(0;0;3)$.
	 \choiceTF
		{\True Một véc-tơ pháp  tuyến của mặt phẳng $(ABC)$ là $\overrightarrow{n}=(3;3;2)$}
		 {\True Mặt phẳng đi qua $C$ và vuông góc với đường thẳng $AB$ có phương trình là $x-y=0$}
		{Mặt phẳng chứa đường thẳng $AB$và vuông góc với $(ABC)$ có  phương trình là $x+y-3z+2=0$}
		{Gọi $M(a;b;c)\in(Oyz)$ sao  cho $\left|\overrightarrow{MA}+\overrightarrow{MB}+\overrightarrow{MC}\right|$ nhỏ nhất, khi đó $3(a+b)+c=5$}
	\loigiai{ 
		\begin{itemchoice}
			\itemch Phương trình mặt phẳng $(ABC)\colon \dfrac{x}{2}+\dfrac{y}{2}+\dfrac{z}{3}=1\Leftrightarrow 3x+3y+2z-6=0$.\\ 
			Suy ra một véc-tơ pháp tuyến của mặt  phẳng $(ABC)$ là $\overrightarrow{n}=(3;3;2)$.
			\itemch Gọi  $(\alpha)$ là mặt phẳng qua $C$ và vuông góc với $AB$.\\
			Mặt phẳng $(\alpha)$ có véc-tơ  pháp tuyến là $\overrightarrow{AB}=(-2;2;0)=-2(1;-1;0)$.\\
			Phương trình của $(\alpha)$ là $ 1(x-0)-1(y-0)+0(z-3)=0\Leftrightarrow  x-y=0$.
			\itemch Gọi $(\beta)$ là mặt phẳng mặt phẳng chứa đường thẳng $AB$ và vuông góc với $(ABC)$  có một véc-tơ pháp tuyến của $(\beta)$ là $\overrightarrow{n}_{\beta}=\left[\overrightarrow{AB},\overrightarrow{n}\right]=(4;4;-12)=4(1;1;-3)$.\\
			 Phương trình của $(\beta)$ là $1(x-2)+1(y-0)-3(z-0)=0\Leftrightarrow x+y-3z-2=0$.
			\itemch Gọi $G$ là  trọng tâm của tam giác $ABC$.\\
			Ta có $G\left(\dfrac{2}{3};\dfrac{2}{3};1\right)$.\\ 
			Vì  $\left|\overrightarrow{MA}+\overrightarrow{MB}+\overrightarrow{MC}\right|=3\left|\overrightarrow{MG}\right|$.\\
			Suy ra $\left|\overrightarrow{MA}+\overrightarrow{MB}+\overrightarrow{MC}\right|_{\min}\Leftrightarrow{\left|\overrightarrow{MG}\right|_{\min}}$.\\ 
			Mà $M\in Oyz\Rightarrow M$ phải là  hình chiếu của $G$ lên $Oyz\Rightarrow M\left(0;\dfrac{2}{3};1\right)$.\\
			Vậy  $3(a+b)+c=2+1=3$.
		\end{itemchoice}
	 }
\end{ex}

% nguồn: Phan1_Lop_12_Hinh_hoc_Chuong_5
\begin{ex}%[2H5V1-5]%[ma:1fdab5e2]
	Trong  hệ trục tọa độ $Oxyz$, cho hai mặt phẳng $(P)\colon 2x-2y+z+1=0$, $(Q)\colon 2x-2y+z-5=0$ và các điểm $A(0;1;1)$, $B(2;0;1)$.
	\choiceTF 
		{\True Hai  mặt phẳng $(P)$ và $(Q)$ song song với nhau}
		{Khoảng  cách $\mathrm{d}\Big(A,(Q)\Big)=4$}
		 {Khoảng cách $\mathrm{d}\Big(P,(Q)\Big)=6$}
		{Cho biết  điểm $C\in(P)$ và đường thẳng $BC$ tạo với mặt phẳng $(P)$ góc $30^\circ$. Khi đó ta có khoảng cách $BC=\dfrac{4\sqrt{3}}{3}$}
	\loigiai{ 
		\begin{itemchoice}
			\itemch Ta có $\overrightarrow{n}_1=(2;-2;1)$ và $\overrightarrow{n}_2=(2;-2;1)$ là các véc-tơ  pháp tuyến của hai mặt phẳng $(P)$ và $(Q)$ tương ứng.\\
			Rõ ràng chúng cùng phương nhau.\\ 
			Mặt khác ta có $1\ne-5$ là các hệ số tự do của phương trình tổng quát  của $2$ mặt phẳng tương ứng.\\
			Từ đó ta có hai mặt phẳng  $(P)$ và $(Q)$ song song nhau.
			\itemch Ta  có $\mathrm{d}\Big(A,(Q)\Big)=\dfrac{\left|0-2+1-5\right|}{\sqrt{4+4+1}}=2$.
			\itemch Ta thấy $A(0;1;1)$ thuộc $(P)$ vì tọa độ điểm $A$ thỏa mãn phương trình mặt phẳng  $(P)$ ta được $20-2\cdot 1+1\cdot 1+1=0$.\\
			Vậy  có $(P)\parallel (Q)$ và $A\in(P)$ nên $\mathrm{d}\Big(P,(Q)\Big)=\mathrm{d}(A,(Q)\Big)=2$.
			\itemch Ta thấy $B(2;0;1)$ thuộc $(Q)$ vì tọa độ điểm $B$ thỏa  mãn phương trình mặt phẳng $(Q)$ ta đucợ $2\cdot 2-2\cdot 0+1\cdot 1-5=0$.\\
			 Gọi $H$ là hình chiếu vuông góc của $B$ lên $(P)$, ta có $BH=\mathrm{d}\\Big
			 (B,(P)\Big)=\mathrm{d}\Big(A,(Q)\Big)=2$.\\
			Vì góc giữa đường thẳng $BC$ và mặt phẳng $(P)$ bằng $30^\circ$ nên suy ra góc  $\widehat{BCH}=30^\circ$.\\
			Do $C\in(P)$ nên xét tam  giác $BHC$ vuông tại $H$ có $BH=2$, $\widehat{BCH}=30^\circ$.\\
			Khi đó ta có  $\sin 30^\circ=\dfrac{BH}{BC}\Rightarrow BC=\dfrac{BH}{\sin 30^\circ}=\dfrac{2}{0{,}5}=4$.
		\end{itemchoice}
	 }
\end{ex}

% nguồn: Phan1_Lop_12_Hinh_hoc_Chuong_5
\begin{ex}%[2H5V1-5]%[ma:9677be3d]
	Trong không gian với hệ trục tọa độ $Oxyz$,  cho hai mặt phẳng $(P)\colon x-3y+2z-1=0$ và $(Q)\colon -2x+y+4z-3=0$.
	 \choiceTF
		{Mặt phẳng $(P)$ song  song mặt phẳng $(Q)$}
		{\True Mặt phẳng $(R)$ có phương trình $14x+8y+5z-2=0$ cùng vuông góc với hai mặt phẳng $(P)$ và  $(Q)$}
		{Mặt phẳng $(\alpha)\colon 2x-6y+4z-9=0$ song song mặt phẳng $(P)\colon x-3y+2z-1=0$ và cách mặt phẳng  $(P)$ một khoảng bằng $1$}
		{\True Mặt  phẳng $(\beta)$ qua hai điểm $A(1;0;1)$, $B(2;-2;1)$ và vuông góc với mặt phẳng $(Q)$ có một véc-tơ pháp tuyến là $\overrightarrow{n}_{\beta}=(8;4;3)$}
	 \loigiai{
		\begin{itemchoice}
			\itemch véc-tơ pháp tuyến của mặt  phẳng $(P)$ là $\overrightarrow{n}_P=(1;-3;2)$.\\
			véc-tơ pháp tuyến  của mặt phẳng $(Q)$ là $\overrightarrow{n}_Q=(-2;1;4)$.\\
			Vì hai véc-tơ pháp tuyến $\overrightarrow{n}_P=(1;-3;2)$ và  $\overrightarrow{n}_Q=(-2;1;4)$ của hai mặt phẳng $(P)$ và $(Q)$ không cùng phương.\\
			Do đó mặt phẳng $(P)$ không song  song mặt phẳng $(Q)$.
			\itemch Mặt phẳng  $(R)$ cùng vuông góc với hai mặt phẳng $(P)$ và $(Q)$ nên có cặp vecto chỉ phương là $\overrightarrow{n}_P=(1;-3;2)$ và $\overrightarrow{n}_Q=(-2;1;4)$.\\
			Mặt phẳng $(R)$ có một véc-tơ pháp tuyến là $\left[\overrightarrow{n}_P,\overrightarrow{n}_Q\right]=(-14;-8;-5)$ hay chọn $\overrightarrow{n}_R=(14;8;5)$.\\ 
			Vậy mặt phẳng $(R)$ có phương trình $14x+8y+5z-2=0$ cùng vuông góc  với hai mặt phẳng $(P)$ và $(Q)$.
			\itemch véc-tơ pháp tuyến  của mặt phẳng $(P)$ là $\overrightarrow{n}_P=(1;-3;2)$.\\
			véc-tơ pháp tuyến của mặt  phẳng $(\alpha)$ là $\overrightarrow{n}_{\alpha}=(2;-6;4)$.\\
			 Ta có $\heva{&\overrightarrow{n}_{\alpha}=2\overrightarrow{n}_P\\&-9\ne 2\cdot (-1)}$ nên mặt phẳng $(\alpha)$ song song mặt phẳng $(P)$.\\
			Lấy điểm $M(1;0;0)$ thuộc mặt phẳng  $(P)\colon x-3y+2z-1=0$.\\
			Lại có $\mathrm{d}\Big(M,(\alpha)\Big)=\dfrac{\left|2\cdot  1-9\right|}{\sqrt{2^2+(-6)^2+4^2}}=\dfrac{\sqrt{14}}{4}$.\\
			Nên mặt phẳng $(\alpha)$ cách mặt phẳng  $(P)$ một khoảng bằng $\dfrac{\sqrt{14}}{4}$.
			\itemch Ta có $\overrightarrow{AB}=(1;-2;0)$ và véc-tơ pháp tuyến  của mặt phẳng $(Q)$ là $\overrightarrow{n}_Q=(-2;1;4)$.\\
			Mặt phẳng $(\beta)$ qua hai điểm $A(1;0;1)$, $B(2;-2;1)$ và vuông góc với mặt phẳng $(Q)$ nên có  cặp vec tơ chỉ phương là $\overrightarrow{AB}$ và $\overrightarrow{n}_Q$.\\
			Do đó mặt phẳng $(\beta)$ có một  véc-tơ pháp tuyến là $\overrightarrow{n}_{\beta}=\left[\overrightarrow{n}_Q,\overrightarrow{AB}\right]=(8;4;3)$.
		\end{itemchoice}
	} 
\end{ex}


%%%%%%%%%%%%%%%% TRẢ LỜI NGẮN %%%%%%%%%%%%%%%%

% nguồn: Phan2_Lop_12_Hinh_hoc_Chuong_5
\begin{bt}%[2H5N1-5]%[ma:54480702]
	Cho mặt  cầu $(S)$ có phương trình $(x+1)^2+(y-2)^2+z^2=16$. Tính khoảng cách từ tâm mặt cầu $(S)$ đến mặt phẳng $(P)\colon 2 x+y-2 z-6=0$.
	 \shortans{2}
	\loigiai{ 
		Mặt cầu  $(S)$ có tâm $I(-1;2;0)$.\\
		Khoảng cách  từ tâm $I$ đến mặt phẳng $(P)$ là $$\mathrm{\,d}(I,(P))=\dfrac{|2\cdot (-1)+2-2\cdot 0-6|}{\sqrt{2^2+1^2+(-2)^2}}=2.$$
	 }
\end{bt}

% nguồn: Phan2_Lop_12_Hinh_hoc_Chuong_5
\begin{bt}%[2H5H1-5]%[ma:127013d3]
	Trong không gian $Oxyz$, cho điểm $A(-2; 6; 0)$ và mặt phẳng $(\alpha)\colon3x+4y+89=0$. Đường thẳng $d$ thay đổi nằm trên mặt phẳng $(Oxy)$ và luôn đi qua điểm $A$. Gọi $H$ là hình chiếu vuông góc của  $M(4;-2; 3)$ trên đường thẳng $d$. Khoảng cách nhỏ nhất từ $H$ đến mặt phẳng $(\alpha)$ bằng bao nhiêu?
	 \shortans{15}
	 \loigiai{
		\begin{center}
			\begin{tikzpicture}[>=stealth,line join=round,line cap=round,font=\footnotesize]
					\path 
						(0,0) coordinate  (I)
						(60:2cm and 1.5cm) coordinate (K) 
						($(K)!2!(I)$) coordinate  (A)
						(-20:2cm and 1.5cm) coordinate  (H)
						(0:2cm and 1.5cm) coordinate (H_0) 
						(K)++(90:2) coordinate (M) 
						 (H_0)++(0:1.1) coordinate (m)
						(H)++(0:1.5) coordinate (n) 
						($(m)!4.5!(n)$)  coordinate (x)
						($(n)!6!(m)$) coordinate (y) 
						($(m)!4!(n)$)  coordinate (X)
						($(n)!5!(m)$) coordinate (Y); 
					 \draw
						(I) ellipse (2cm and  1.5cm)
						 (A)--(K)--(H)--cycle
						(I)--(m) 
						 (H)--(n)
						 (x)--(y)
						 (K)--(M)--(H)
						(X)--++(90:2)--($(Y)+(90:2)$)--(Y) 
						(X)--++(180:7) 
						($(X)+(160:1)$) node{$(Oxy)$} 
						($(Y)+(75:1)$) node{$(\alpha)$}; 
					\foreach  \m/\n in{I/90,K/200,H/-60,A/-130,H_0/40,M/90}
					 \fill (\m) circle (1pt)+(\n:3mm) node{$\m$};
				\end{tikzpicture}
		\end{center}
		Gọi $K$ là  hình chiếu vuông góc của $M$ lên $(Oxy)$, suy ra $K(4;-2; 0)$.\\
		 Vì $\heva{&AH\perp MK\left(\mathrm{do } MK\perp (Oxy)\right)\\&AH\perp MH} \Rightarrow AH\perp (MKH)\Rightarrow AH\perp KH$.\\
		 Khi đó $H$ luôn thuộc đường tròn $(C)$ có tâm là trung điểm $I(1; 2; 0)$ của đoạn $AK$, bán kính $R=\dfrac{AK}{2}=5$.\\
		 Gọi $\Delta=(\alpha)\cap (Oxy)$, ta thấy $(\alpha)\perp (Oxy)$ (vì $\overrightarrow{n}_{(\alpha)} \cdot \vec{k}=0$).\\
		Khi đó $\mathrm{d}(I, \Delta)=\mathrm{d}\left(I,  (\alpha)\right)=20$.\\
		Khoảng cách nhỏ nhất từ $H$ đến mặt phẳng $(\alpha)$ là 
		 \[\mathrm{d}\left(H, (\alpha)\right)_{min}=\mathrm{d}(H, \Delta)_{min}=\mathrm{d}\left(H_0, \Delta \right)=\mathrm{d}(I, \Delta)-R=20-5=15.\]
	 }
\end{bt}

% nguồn: Phan2_Lop_12_Hinh_hoc_Chuong_5
\begin{bt}%[2H5H1-5]%[ma:d8e7b294]
	Trong không gian $Oxyz$,  cho điểm $M(-1; 2;-2)$. Tính tổng khoảng cách từ điểm $M$ đến các mặt phẳng tọa độ.
	\shortans{5} 
	\loigiai{ 
		Ta  gọi $M_1$, $M_2$, $M_3$ lần lượt là hình chiếu của $M(-1;2;-2)$ lên mặt phẳng $(Oxy)$, $(Oyz)$ và $(Oxz)$.\\
		 Khi đó $M_1(-1;2;0)$, $M_2(0;2;-2)$ và $M_3(-1;0;-2)$.\\
		Ta có 
		 \allowdisplaybreaks
		\begin{eqnarray*}
			\heva{ 
			&\mathrm{d}\big(M,(Oxy)\big)=MM_1=2\\ 
			 &\mathrm{d}\big(M,(Oyz)\big)=MM_2=1\\
			&\mathrm{d}\big(M,(Oxz)\big)=MM_3=2. 
			} 
		\end{eqnarray*}
		Vậy tổng khoảng cách từ điểm  $M$ đến các mặt phẳng tọa độ bằng $2+1+2=5$.
	} 
\end{bt}

% nguồn: Phan2_Lop_12_Hinh_hoc_Chuong_5
\begin{bt}%[2H5H1-5]%[ma:09b0ef1b]
	Bản vẽ thiết kế của một công trình được vẽ trong hệ trục  tọa độ $Oxyz$. Sàn nhà thuộc mặt phẳng $Oxy$, đường ống nước thẳng đi qua hai điểm $A(1; 2;-1)$ và $B(5; 6;-2)$. Trên đường ống nước gắn một cái khóa tại trung điểm của $AB$. Tính khoảng cách từ cái khóa đến mặt sàn nhà
	\shortans{1,5} 
	 \loigiai{
		Gọi $I$ là trung điểm của $AB$, khi đó $I\left(3;4;-\dfrac{3}{2}\right)$.\\ 
		Suy ra tọa độ của cái  khóa là $I\left(3;4;-\dfrac{3}{2}\right)$.\\
		Khoảng cách từ  cái khóa đến mặt sàn nhà là khoảng cách từ điểm $I$ đến mặt phẳng $Oxy$. Gọi $H(3;4;0)$ là hình chiếu của điểm $I$ lên mặt phẳng $Oxy$, khi đó
		\[ 
		 \mathrm{d}\left(I,(Oxy)\right)=IH=\dfrac{3}{2}=1{,}5.
		\] 
	 }
\end{bt}

% nguồn: Phan2_Lop_12_Hinh_hoc_Chuong_5
\begin{bt}%[2H5H1-5]%[ma:60c6976d]
	Biết góc quan sát ngang của một camera là $116^\circ$. Trong  không gian $Oxyz$, camera được đặt tại điểm $A(2;1;5)$ và chiếu thẳng về phía mặt phẳng $(P)\colon 2x-y-2z+13=0$.
	\begin{center}
		\begin{tikzpicture}[line cap=round,line join=round,>=stealth,x=1.0cm,y=1.0cm,scale=0.7]
				 \def\a{2}
				 \def\b{1}
				 \def\h{4}
				 \pgfmathsetmacro\g{asin (\b/\h)}
				\pgfmathsetmacro\xo{\a  *cos (\g)}
				\pgfmathsetmacro\yo{\b  *sin (\g)}
				 \draw[dashed]
					(\xo,\yo) arc  (\g:180-\g:{\a} and
					{\b}) 
					 (180:\a) coordinate (B) node[left]{$ $}
					 --(0:\a) coordinate (A)
					(90:\h)  coordinate (S)
					--(0:0) coordinate  (O)
					($(B)+(-2,-1.3)$)  coordinate (E)
					($(E)+(60:3cm)$) coordinate (F) 
					($(F)+(0:6cm)$) coordinate  (G)
					($(E)+(0:6cm)$) coordinate (H) 
					(intersection of F--G and  S--B) coordinate (M)
					(intersection of F--G and S--A) coordinate (N); 
				\draw  (90:\h)--(-\xo,\yo) arc
					 (180-\g:360+\g:{\a} and {\b})--cycle;
				\foreach \x/\g in{}  \fill (\x) circle (1pt)
					($(\x)+(\g:3mm)$) node{\small $\x$}; 
				\draw (M)--(F)--(E)--(H)--(G)--(N); 
				\draw  [dashed] (M)--(N);
				\draw pic[draw,angle radius=6mm,angle eccentricity=0.62,"$P$"] {angle  = H--E--F};
				\draw pic[draw,angle radius=12mm,angle eccentricity=0.62,"$116^\circ$"] {angle =  B--S--A};
				\draw (S) node [above]  {Camera};
			\end{tikzpicture}
	\end{center}
	Hỏi vùng quan sát được trên mặt phẳng $(P)$ của camera là hình tròn có diện tích bằng  bao nhiêu? (Làm tròn kết quả đến chữ số hàng phần chục).
	 \shortans{32,2}
	\loigiai{ 
		Gọi $H$ là hình chiếu vuông góc  của $A$ lên mặt phẳng $(P)$ và $B,C$ là hai điểm biên của vùng quan sát sao cho $\widehat{BAC}=116^\circ$.
		Vì camera quan sát đối xứng  sang hai phía nên
		\[ 
		 \widehat{BAH}=\widehat{CAH}=\dfrac{116^\circ}{2}=58^\circ.
		\] 
		Trước  hết, ta tính khoảng cách từ $A$ đến mặt phẳng $(P)$:
		 \[
		AH=d\bigl(A,(P)\bigr) 
		 =\dfrac{|2\cdot2-1-2\cdot5+13|}{\sqrt{2^2+(-1)^2+(-2)^2}}
		 =\dfrac{|4-1-10+13|}{\sqrt{4+1+4}}
		=\dfrac{6}{3}=2\  \text{ (đvđd)}.
		 \]
		Xét tam giác vuông $ABH$ tại  $H$, ta có
		\[ 
		 \tan\widehat{BAH}=\dfrac{BH}{AH}
		\Rightarrow  BH=AH\tan58^\circ=2\tan58^\circ.
		\] 
		Suy ra bán kính của hình tròn vùng quan sát  là
		\[ 
		 r=BH=2\tan58^\circ.
		 \]
		Diện tích vùng quan  sát trên mặt phẳng $(P)$ là
		 \[
		 S=\pi r^2=\pi\left(2\tan58^\circ\right)^2\approx32{,}2\ \text{ (đvdt)}.
		\] 
		Vậy diện tích vùng quan sát của camera trên mặt phẳng $(P)$  xấp xỉ $32{,}2$.
	} 
\end{bt}

% nguồn: Phan2_Lop_12_Hinh_hoc_Chuong_5
\begin{bt}%[2H5H1-5]%[ma:6559cd33]
	Trong không gian $Oxyz$, cho điểm $M(1;-1;0)$ và mặt phẳng $(P)\colon x-2y+2z-9=0$. Tính khoảng cách từ điểm $M$ đến mặt  phẳng $(P)$.\shortans{2}
	\loigiai{ 
		Khoảng cách từ $M(1; -1; 0)$ đến mặt phẳng $(P)\colon x - 2y + 2z - 9 =  0$ là
		$$\mathrm{d}(M, (P)) =  \dfrac{|1\cdot1 - 2\cdot(-1) + 2\cdot0 - 9|}{\sqrt{1^2 + (-2)^2 + 2^2}}= 2.$$
	} 
\end{bt}

% nguồn: Phan2_Lop_12_Hinh_hoc_Chuong_5
\begin{bt}%[2H5V1-5]%[ma:c7e26a4b]
	Trong không gian $Oxyz$, cho điểm $M(1;4;9)$. Gọi $(P)$ là mặt phẳng đi qua $M$ và cắt ba tia $Ox$, $Oy$, $Oz$ lần lượt tại các điểm $A$, $B$, $C$ (khác $O$) sao cho $OA+OB+OC$ đạt giá trị nhỏ nhất. Tính khoảng cách $d$ từ gốc tọa độ $O$ đến mặt phẳng $(P)$ (kết quả là tròn đến hàng  phần trăm).
	\shortans{5,14} 
	\loigiai{ 
		Giả sử $A(a;0;0)$, $B(0;b;0)$, $C(0;0;c)$ với $a,b,c  > 0$.\\
		Phương trình mặt phẳng $(P)\colon \dfrac{x}{a}  + \dfrac{y}{b} + \dfrac{z}{c} = 1$; $OA+OB+OC = a+b+c$.\\
		Ta có $M(1;4;9) \in (P) \Rightarrow \dfrac{1}{a} + \dfrac{4}{b} + \dfrac{9}{c} = 1$. Áp dụng BĐT  Bunhiacopxki
		\begin{align*}
			\left( \dfrac{1}{a} + \dfrac{4}{b}  + \dfrac{9}{c}\right)(a+b+c) &= \left[\left(\sqrt{\dfrac{1}{a}}\right)^2 + \left(\sqrt{\dfrac{4}{b}}\right)^2 + \left(\sqrt{\dfrac{9}{c}}\right)^2\right]\left((\sqrt{a})^2+(\sqrt{b})^2+(\sqrt{c})^2\right)\\&\ge (1+2+3)^2\Rightarrow a+b+c \ge 36.
		\end{align*}
		Dấu \lq \lq $ = $ \rq \rq  xảy ra khi
		\[ 
		\heva{&  \dfrac{1}{a} + \dfrac{4}{b} + \dfrac{9}{c} = 1 \\& \dfrac{1}{a} = \dfrac{2}{b} = \dfrac{3}{c}}
		 \xrightarrow{a+b+c=49}
		\heva{&a=6 \\&b=12 \\&c=18}\] 
		Nên $ (P)\colon  \dfrac{x}{6} + \dfrac{y}{12} + \dfrac{z}{18} = 1$. \\
		Vậy $  d = \dfrac{36}{7} \approx 5{,}14.$
	 }
\end{bt}

% nguồn: Phan2_Lop_12_Hinh_hoc_Chuong_5
\begin{bt}%[2H5V1-5]%[ma:ee981501]
	Trong không gian $Oxyz$, cho điểm $M(1;2;3)$. Mặt phẳng $(P)\colon x + Ay + Bz  + C = 0$ chứa trục $Oz$ và cách điểm $M$ một khoảng lớn nhất, khi đó tính tổng $A+B+C$.
	\shortans{2} 
	\loigiai{ 
		Vì $(P)$ chứa trục $Oz$  nên luôn có $d(M;(P)) \le d(M;Oz)$.\\
		Suy ra $d(M;(P))$ đạt  giá trị lớn nhất bằng $d(M;Oz)=MH$, với $H$ là hình chiếu của $M$ trên trục $Oz$.\\
		Ta có $H(0;0;3)$. Vậy $(P)$ đi qua $H(0;0;3)$, có véc tơ  pháp tuyến $\overline{MH}(-1;-2;0)$.
		$(P)\colon -x-2y=0 \Leftrightarrow x+2y=0 \Rightarrow A=2; B=C=0 \Rightarrow A+B+C=2$. 
	 }
\end{bt}

% nguồn: Phan2_Lop_12_Hinh_hoc_Chuong_5
\begin{bt}%[2H5V1-5]%[ma:c89fd792]
	Trong không gian với hệ trục tọa độ $Oxyz$, cho hai điểm $A(2;3;4)$, $B(6;-3;0)$ và mặt phẳng $(P)$ sao cho $\mathrm{d}(B;(P)) = 2\mathrm{d}(A;(P))$, $(P)$ cắt $AB$ tại $I(a;b;c)$ nằm trên  đoạn thẳng $AB$. Tính $S = a+b+c$.
	 \shortans{7}
	 \loigiai{
		Gọi $H, K$ lần lượt là hình chiếu vuông góc của $A$ và $B$ lên  mặt phẳng $(P)$.\\
		Suy ra $BK = d(B;(P))$, $AH = d(A;(P))$  do $AH \perp (P); BK \perp (P)$ nên $AH \parallel BK$.\\
		 Suy ra $A, H, K, B$ đồng phẳng và ta có $\triangle IKB \sim \triangle IHA$ (góc - góc).\\
		Suy ra $ \dfrac{BI}{AI} = \dfrac{BK}{AH} = \dfrac{d(B;(P))}{d(A;(P))} = 2 \Rightarrow BI = 2AI$ mà $I$ thuộc đoạn $AB$ nên $\overrightarrow{BI}  = -2\overrightarrow{AI}$.\\
		Ta có $\overrightarrow{BI}(a-6; b+3;  c)$, $\overrightarrow{AI}(a-2; b-3; c-4)$.\\
		$$\overrightarrow{BI} = -2\overrightarrow{AI} \Leftrightarrow 
		\heva{& a-6 = -2(a-2) \\&b+3 =  -2(b-3) \\& c = -2(c-4)}
		 \Leftrightarrow
		\heva{& a = \dfrac{10}{3}  \\&b=1 \\&c = \dfrac{8}{3}}\Rightarrow a+b+c=7.$$
		Vậy $S=7$. 
	 }
\end{bt}

% nguồn: Phan2_Lop_12_Hinh_hoc_Chuong_5
\begin{bt}%[2H5V1-5]%[ma:053e60de]
	Trong không  gian $Oxyz$, cho hình chóp $S.ABCD$ có đáy $ABCD$ là hình thang vuông tại $A(0;0;0)$, $B(6;0;0)$ và cạnh bên $SA$ vuông góc với mặt phẳng đáy. Biết rằng $D(0;x;0)$, với $x>0$ thỏa mãn $AD=2AB=2BC$, góc giữa đường thẳng $SC$ và mặt phẳng đáy bằng $45^\circ$. Tính khoảng cách từ $B$ đến mặt phẳng $(SCD)$.
	 \shortans{3}
	 \loigiai{
		\begin{center}
			\begin{tikzpicture}[scale=1,font=\footnotesize,line join=round,line cap=round,>=stealth]
					 \path
						(0,0)  coordinate (A) node[above left]{$A$}
						(-2,-1)  coordinate (B) node[above left]{$B$}
						(4,0) coordinate (M) node[above]{$M$} 
						 ($(B)+(M)-(A)$) coordinate (C) node[below]{$C$}
						(A)+(0,4) coordinate  (S) node[above right]{$S$}
						(8,0) coordinate  (D) node[below]{$D$};
					\draw (S)--(D)--(C)--cycle (S)--(B)--(C) 
						(S)--(D); 
					\draw[dashed] (A)--(S) 
						 (B)--(A)--(D)
						 (A)--(C)--(M);
					\draw[->] (S)--($(A)!1.2!(S)$) node[above left]{$z$}; 
					\draw[->]  (B)--($(A)!1.4!(B)$) node[above left]{$x$};
					\draw[->]  (D)--($(A)!1.2!(D)$) node[below right]{$y$};
					 \foreach \x in{A,B,C,D,S,M}\draw (\x) circle (1pt);
				\end{tikzpicture}
		\end{center}
		Gọi $M$ là  trung điểm của $AD$ thì $\vv{AM}=\vv{BC}$, suy ra $ABCM$ là hình vuông.\\ Do đó $M(0;6;0)$, $C(6;6;0)$ và $D(0;12;0)$.\\
		Góc  giữa đường thẳng $SC$ và mặt phẳng đáy bằng $45^\circ$, suy ra $\widehat{SCA}=45^\circ$\\
		$\Rightarrow  \triangle ASC$ vuông cân tại $A$ nên $SA=AC=AB\sqrt{2}=6\sqrt{2}$.\\
		Vậy  $S(0;0;6\sqrt{2})$.\\
		 Giả sử phương trình mặt phẳng $(SCD)$ là $ax+by+cz+d=0$.\\
		 Ta có
		$\heva{&c\cdot 6\sqrt{2}-12=0\\&  6a+6b-12=0\\&12b-12=0}\Leftrightarrow\heva{&a=1\\&b=1\\&c=\sqrt{2}.}$\\
		Do đó, phương trình mặt  phẳng $(SCD)$ là $x+y+\sqrt{2}z=12$.\\
		Khoảng cách từ $B$ đến mặt phẳng $(SCD)$ là $\mathrm{d}\left(B,(SCD) \right)=\dfrac{\left|6+0+\sqrt{2\cdot 0 -12}\right|}{2}=3$. 
	} 
\end{bt}

% nguồn: Phan2_Lop_12_Hinh_hoc_Chuong_5
\begin{bt}%[2H5V1-5]%[ma:66b763cd]
	Trong không gian $Oxyz$, cho mặt phẳng $(\alpha)\colon 2x + 3y + z + 1 = 0$. Gọi $(P)$ là mặt phẳng song song với $(\alpha)$, cắt các tia $Ox, Oy, Oz$ lần lượt tại các điểm $A$, $B$,  $C$ sao cho thể tích khối tứ diện $OABC$ bằng $6$. Tính khoảng cách từ gốc tọa độ $O$ đến mặt phẳng $(P)$ (kết quả làm tròn đến hàng phần trăm).
	\shortans{1,60} 
	\loigiai{ 
		Do $(P)$ song song với $(\alpha)$ nên $(P)$ có  phương trình $2x+3y+z+m=0$, điều kiện $m \neq 1$.\\
		Khi đó $(P)$ cắt các tia $Ox$, $Oy$, $Oz$ lần lượt tại các điểm là $A\left(- \dfrac{m}{2};0;0\right)$, $B\left(0;- \dfrac{m}{3};0\right)$, $C(0;0;-m)$, với  $m < 0$.\\
		Thể tích khối tứ  diện $OABC$ bằng $6$ nên $ \dfrac{1}{6}OA\cdot OB \cdot OC = 6$.
		\begin{align}
			&\Leftrightarrow \dfrac{1}{6} \cdot \left|- \dfrac{m}{2}\right| \cdot \left|- \dfrac{m}{3}\right| \cdot |-m| = 6 \Leftrightarrow - \dfrac{m^3}{36} = 6 \text{ (do } m<0)\\&\Leftrightarrow m^3 = -216 \Leftrightarrow m=-6 \text{  (thỏa mãn)}
		\end{align}
		Ta có $(P)\colon 2x+3y+z-6=0 \Rightarrow d(O;(P))  = \dfrac{|2\cdot 0+3\cdot 0+0-6|}{\sqrt{2^2+3^2+1^2}} = \dfrac{6}{\sqrt{14}} \approx 1{,}60.$
	} 
\end{bt}

% nguồn: Phan2_Lop_12_Hinh_hoc_Chuong_5
\begin{bt}%[2H5C1-5]%[ma:1ffb60dc]
	Trong không gian $Oxyz$, cho các điểm $A(1;-2;-1)$, $B(-3;-1;2)$, $C(3;2;1)$ và điểm $D(-2;1;3)$. Gọi $M$ là điểm thay đổi trên mặt phẳng $(Oxz)$. Giá trị lớn nhất của  biểu thức $P=MA^2-2MB^2-MC^2+MD$ bằng bao nhiêu?
	\shortans{-9} 
	\loigiai{ 
		Tìm tọa độ điểm  $I$ sao cho $\overrightarrow{IA}-2\overrightarrow{IB}-\overrightarrow{IC}=\overrightarrow{O}$\\
		Ta  có $\overrightarrow{IA}-2\overrightarrow{IB}-\overrightarrow{IC}=\overrightarrow{O}\Leftrightarrow \overrightarrow{OA}-2\overrightarrow{OB}-\overrightarrow{OC}=-2\overrightarrow{OI}\Leftrightarrow\overrightarrow{OI}=\dfrac{-1}{2}\overrightarrow{OA}+\overrightarrow{OB}+\dfrac{1}{2}\overrightarrow{OC}$.\\
		Suy ra $\heva{&x_I=-\dfrac{1}{2}x_A+x_B+\dfrac{1}{2}x_c  \\
		 &y_I=-\dfrac{1}{2}y_A+y_B+\dfrac{1}{2}y_c \\
		&z_I=-\dfrac{1}{2}z_A+z_B+\dfrac{1}{2}z_c}\Leftrightarrow\heva{&x_I=-\dfrac{1}{2}+(-3)+\dfrac{3}{2}\\ 
		&y_I=1+(-1)+1\\ 
		&z_I=\dfrac{1}{2}+2+\dfrac{1}{2}}\Leftrightarrow\heva{&x_I=-2\\ 
		 &y_I=1\\
		&z_I=3.}$\\ 
		Do  đó $I(-2;1;3)$.\\
		Hơn  nữa
		\begin{eqnarray*}
			 P&=&MA^2-2MB^2-MC^2+MD={\overrightarrow{MA}}^2-2{\overrightarrow{MB}}^2-{\overrightarrow{MC}}^2+\overrightarrow{MD}\\
			&=&\left(\overrightarrow{IA}-\overrightarrow{IM}\right)^2-2\left(\overrightarrow{IB}-\overrightarrow{IM}\right)^2-\left(\overrightarrow{IC}-\overrightarrow{IM}\right)^2+\overrightarrow{MD}\\ 
			 &=&IA^2-2IB^2-IC^2-2IM^2+MD.
		\end{eqnarray*}
		Ta có $IA^2-2IB^2-IC^2=-8$  và $D\equiv M$ nên $P=-8-2IM^2+IM$.\\
		Dặt $t=IM$.\\ 
		Gọi $H(-2;0;3)$ là hình chiếu vuông góc của điểm $I$ trên mặt phẳng $(Oxz)$, ta có $IH=1$.\\ 
		Khi đó, $IM\geq IH$ hay  $t\geq IH\Rightarrow t\geq1$.\\
		 Xét hàm số $f\left(t\right)=-2t^2+t-8,t\in\left[1;+\infty\right)$.\\
		Ta có $\displaystyle\max_{\left[1;+\infty\right)}f\left(t\right)=f\left(1\right)=-9$.\\ 
		Vậy giá trị lớn nhất của biểu thức  $P$ là $-9$.
	 }
\end{bt}


%%%%%%%%%%%%%%%% TỰ LUẬN %%%%%%%%%%%%%%%%

% nguồn: Phan2_Lop_12_Hinh_hoc_Chuong_5
\begin{bt}%[2H5H1-5]%[2H5H1-3]%[ma:2cb0144d]
	Cho bốn điểm $A(-2; 6; 3)$, $B(1; 0;  6)$, $C(0; 2;-1)$, $D(1; 4; 0)$.
	\begin{enumerate}
		\item Viết phương trình  mặt phẳng $(B C D)$. Suy ra $A B C D$ là một tứ diện.
		\item Tính chiều cao $A  H$ của tứ diện $A B C D$.
		\item Viết phương trình mặt  phẳng $(\alpha)$ chứa $A B$ và song song với $C D$.
	\end{enumerate}
	 \loigiai{
		\begin{enumerate}
			\item  $\overrightarrow{BC}=\left(-1;2;-7\right)$, $\overrightarrow{BD}=\left(0;4;-6\right)$;
			$\left[\overrightarrow{BC},\overrightarrow{BD}  \right]=\left(16; -6; -4\right)$.\\
			Mặt phẳng $(B C D)$ đi qua $B\left(1; 0; 6\right)$  và nhận $\left[\overrightarrow{BC},\overrightarrow{BD} \right]=\left(16; -6; -4\right)$ làm vectơ pháp tuyến.\\
			Phương  trình mặt phẳng $\left(BCD\right)$ là $$16\left(x-1\right)-6\left(y-0\right)-4\left(z-6\right)=0\Leftrightarrow 8x-3y-2z+10=0.$$
			Ta có $-16-18-6+10=-30\ne 0$ nên $A\left(-2; 6; 3\right)\notin \left(BCD\right)$. Suy ra $A B C D$ là một  tứ diện.
			\item Chiều cao $A H$ của tứ diện $A B C D$ là $$AH=\mathrm{d}\left(A,\left(BCD\right)\right)=\dfrac{\left| -16-18-6+10\right|  }{\sqrt{8^2+(-3)^2+(-2)^2}}=\dfrac{30}{\sqrt{77}}.$$
			\item $\overrightarrow{AB}=\left(3;-6;3\right)$,  $\overrightarrow{CD}=\left(1;2;1\right)$.\\
			Mặt phẳng $(\alpha)$ đi qua $B\left(1; 0; 6\right)$ và nhận $\left[\overrightarrow{AB},\overrightarrow{CD} \right]=\left(-12; 0; 12\right)$ làm vectơ pháp tuyến.\\ Phương trình $\left(\alpha\right)$ là 
			$-12\left(x-1\right)+0\left(y-0\right)+12\left(z-6\right)=0\Leftrightarrow -x+z-5=0$.\\ 
			Vậy phương  trình mặt phẳng $(\alpha)$ chứa $A B$ và song song với $C D$ là $-x+z-5=0$.
		\end{enumerate}
	 }
\end{bt}

% nguồn: Phan2_Lop_12_Hinh_hoc_Chuong_5
\begin{bt}%[2H5H1-5]%[ma:cf8f6fe0]
	\immini 
	{Một chiếc lều hình chóp có đáy là hình vuông, mỗi cạnh dài $200$ cm. Đỉnh lều nằm thẳng đứng phía trên tâm của hình vuông, và chiều cao của lều là $206$ cm. Người ta dùng $4$ cọc bằng nhau nối từ $4$ góc của đáy đến  đỉnh lều để dựng lều. Chiều dài tối thiểu của mỗi cọc là bao nhiêu cm (làm tròn đến hàng đơn vị)?
	}{\begin{tikzpicture}
		 \def\a{5}
		\def\h{4.5} 
		\path (0:0) coordinate  (A)
			++(-30:3) coordinate  (D)
			++(-175:\a) coordinate  (C)
			($(A)+(C)-(D)$) coordinate  (B)
			($(B)!0.3!(C)$)  coordinate (E)
			($(E)+(40:1.5)$) coordinate (K) 
			($(C)!0.3!(B)$) coordinate (F) 
			(intersection of A--C and B--D) coordinate  (O)
			 ($(O)+(90:\h)$) coordinate (S)
			($(C)+(90:1)$) coordinate (H) 
			(-2.85,-.7)  coordinate (T);
		 \draw[teal,dashed] (B)--(A)--(D)
			(A)--(S); 
		\draw[teal] (B)--(E) -- (K)--(F)--(C)--(D) 
			(B)--(S) 
			 (C)--(S)
			(D)--(S); 
		\draw[teal]  (K).. controls (T).. (H);
		\draw[teal] (F).. controls ($(T)+(-90:.4)$).. (H); 
	\end{tikzpicture}
	} 
	\loigiai{ 
		\begin{center}
			\begin{tikzpicture}[scale=1,font=\footnotesize,line join=round,line cap=round,>=stealth]
					 \def\a{5}
					\def\h{4.5} 
					\path (0:0) coordinate (A) 
						++(-30:3) coordinate (D) 
						++(-175:\a) coordinate (C) 
						($(A)+(C)-(D)$)  coordinate (B)
						($(B)!0.3!(C)$)  coordinate (E)
						($(E)+(40:1.5)$)  coordinate (K)
						($(C)!0.3!(B)$) coordinate  (F)
						(intersection of A--C and B--D) coordinate  (O)
						($(O)+(90:\h)$) coordinate  (S)
						($(C)+(90:1)$)  coordinate (H)
						(-2.85,-.7)  coordinate (T);
					 \draw[teal,dashed] (B)--(A)--(D)
						 (O)--(A)--(S);
					 \draw[teal] (B)--(E) -- (K)--(F)--(C)--(D)
						(B)--(S) 
						(C)--(S) 
						(D)--(S); 
					\draw[teal] (K).. controls (T).. (H); 
					\draw[teal] (F).. controls ($(T)+(-90:.4)$)..  (H);
					\foreach \x/\g  in {A/30,B/180,C/225,D/-30,S/90,O/-90} \draw[teal,fill=blue] (\x) circle (1pt)
						($(\x)+(\g:3.75mm)$) node{$\x$}; 
					 \draw[teal,<->] ([shift={(0,-0.25cm)}]C)--([shift={(0,-0.25cm)}]D) node[midway,below]{$200$};
					\draw[teal,dashed,<->] (S)--(O) node[midway,right]{$206$}; 
				\end{tikzpicture}
		\end{center}
		Gọi $O$ là tâm hình  vuông $ABCD$.\\
		Theo đề, ta có  $SO\perp(ABCD)$, $SO=206$, $AB=200$.\\
		Ta  có $AO=\dfrac{1}{2}AC=\dfrac{200\sqrt{2}}{2}=100\sqrt{2}$.\\
		Chiều dài  tối thiểu mỗi cột là $SA=\sqrt{SO^2+OA^2}=\sqrt{206^2+\left(100\sqrt{2}\right)^2}=\sqrt{62\,436}\approx250$ cm.
	 }
\end{bt}

% nguồn: Ok_TuanLuan_Toan11_Lan2
\begin{ex}%[2H5V1-5]%[ma:5fba8f65]
	Cho hình lăng trụ tam giác đều $ABC.DEF$ biết rằng $AB=6$, $AD=2$. Gọi ba điểm $M$, $N$, $P$ lần lượt là trung điểm $DE$, $DF$, $BC$. Lập hệ trục tọa độ $Oxyz$ như  hình bên. Gọi điểm $S$ thỏa mãn hệ thức $\overrightarrow{SA}+2\overrightarrow{SB}+\overrightarrow{SC}=\overrightarrow{0}$. Tính khoảng cách từ điểm $S$ đến mặt phẳng $(MNP)$.
	\loigiai{ 
		 Ta có $PA=\dfrac{AB\sqrt 3}{2}=3\sqrt 3$; $PB=PC=\dfrac{BC}{2}=3$; $BE=BF=AD=2$.\\
		Đặt hệ trục  tọa độ như hình vẽ
		\begin{center}
			\begin{tikzpicture}[line join=round,line cap=round,>=stealth,scale=0.8]
					\coordinate  (P) at (0,0);
					\coordinate  (y) at ($(P)+(-5:4)$);
					\coordinate  (z) at ($(P)+(90:5)$);
					\coordinate (x)  at ($(P)+(-150:3)$);
					 \coordinate (A) at ($(P)+(-5:3)$);
					\coordinate (B) at  ($(P)+(-150:2)$);
					\coordinate (C) at ($(P)+(30:2)$); 
					\coordinate  (D) at ($(A)+(90:3)$);
					 \coordinate (E) at ($(B)+(90:3)$);
					\coordinate (F)  at ($(C)+(90:3)$);
					\coordinate (M)  at ($(D)!0.5!(E)$);
					\coordinate (N)  at ($(D)!0.5!(F)$);
					 \coordinate (P') at ($(E)!0.5!(F)$);
					\draw[->] (B) -- (x)  node[left] {$x$};
					\draw[->] (A)  -- (y) node[below] {$y$};
					\draw[->]  (P') -- (z) node[left] {$z$};
					\draw (D)--(E)--(F)--cycle (B)--(E) 
						 (A)--(D)
						 (B)--(A)
						(M)--(N); 
					\draw [dashed]  (B)--(C)--(A)
						(C)--(F) 
						(P)--(A) 
						 (P)--(P')
						(M)--(P)--(N); 
					\foreach  \i/\g in {A/-90,B/-90,C/30,D/0,E/180,F/90,M/90,N/30,P/180}{\draw[fill](\i) circle (1pt)
						($(\i)+(\g:3mm)$) node[scale=1]{$\i$};} 
				\end{tikzpicture}
		\end{center}
		Khi đó $P(0;0;0)$, $A(3\sqrt{3};0;0)$,  $B(0;3;0)$, $C(0;-3;0)$, $E(0;3;2)$, $F(0;-3;2)$, $D(3\sqrt{3};0;2)$.\\
		Suy ra $M\left(\dfrac{3\sqrt{3}}{2};\dfrac{3}{2};2\right)$, $N\left(\dfrac{3\sqrt{3}}{2};\dfrac{-3}{2};2\right)$.\\ 
		 Ta có $\overrightarrow{PM}=\left(\dfrac{3\sqrt{3}}{2};\dfrac{3}{2};2\right)$; $\overrightarrow{PN}=\left(\dfrac{3\sqrt{3}}{2};\dfrac{-3}{2};2\right)$.\\
		Suy ra  $\left[\overrightarrow{PM},\overrightarrow{PN}\right]=\left(6;0;\dfrac{-9\sqrt{3}}{2}\right)\Rightarrow \overrightarrow{n}=(4;0;-3\sqrt{3})$.\\
		Phương trình mặt phẳng  $(MNP)$ là $4x-3\sqrt{3}z=0$.\\
		Gọi $I$, $J$ lần lượt là trung điểm cạnh $AC$, $IB\Rightarrow I\left(\dfrac{3\sqrt{3}}{2};\dfrac{-3}{2};0\right)$,  $J\left(\dfrac{3\sqrt{3}}{4};\dfrac{3}{4};0\right)$.\\
		Ta có $\overrightarrow{SA}+2\overrightarrow{SB}+\overrightarrow{SC}=\overrightarrow{0}\Leftrightarrow 2\overrightarrow{SI}+2\overrightarrow{SB}=\overrightarrow{0}\Leftrightarrow 4\overrightarrow{SJ}=\overrightarrow{0}\Leftrightarrow S\equiv  J$.\\
		 Do đó $S\left(\dfrac{3\sqrt{3}}{4};\dfrac{3}{4};0\right)$.\\
		 Vậy khoảng cách từ điểm $S$ đến mặt phẳng $(MNP)$ là
		\[\mathrm{d}\Big(S,(MNP)\Big)=\dfrac{\left|4\cdot  \dfrac{3\sqrt 3}{4}-3\sqrt 3\cdot 0\right|}{\sqrt{4^2+\left(-3\sqrt 3\right)^2}}=\dfrac{3\sqrt{129}}{43}.\]
	 }
\end{ex}

% nguồn: Ok_TuanLuan_Toan11_Lan2
\begin{ex}%[2H5V1-5]%[ma:c39ea26d]
	Cho hình lập phương $ABCD.EFGH$ cạnh $1$. Điểm $M$ được cho thỏa mãn hệ thức véc-tơ $\overrightarrow{AM}+\overrightarrow{AE}=3\overrightarrow{CD}$. Tính khoảng cách từ điểm $M$ đến mặt phẳng $(EBD)$. 
	\loigiai{ 
		\begin{center}
			\begin{tikzpicture}[line join=round,line cap=round,>=stealth,scale=0.8]
					 \coordinate (A) at (0,0);
					\coordinate (y) at ($(A)+(0:5)$); 
					\coordinate (z)  at ($(A)+(90:5)$);
					\coordinate (x)  at ($(A)+(-120:3)$);
					\coordinate (B) at ($(A)+(-120:1.5)$); 
					\coordinate (D) at ($(A)+(0:4)$); 
					\coordinate (C) at  ($(B)+(D)-(A)$);
					\coordinate (E)  at ($(A)+(90:4)$);
					 \coordinate (F) at ($(B)+(90:4)$);
					\coordinate (G) at  ($(C)+(90:4)$);
					\coordinate  (H) at ($(D)+(90:4)$);
					 \fill[gray!20](B)--(D)--(E);
					\draw[->] (B) --  (x) node[left] {$x$};
					\draw[->] (D) -- (y)  node[below] {$y$};
					\draw[->] (E) --  (z) node[left] {$z$};
					\draw  (B)--(C)--(D)
						 (E)--(F)--(G)--(H)--cycle (B)--(F)
						 (C)--(G)
						 (D)--(H);
					 \draw [dashed](A)--(B)
						 (A)--(D)
						 (A)--(E)
						(B)--(D)--(E)--cycle; 
					\foreach  \i/\g in {A/45,B/180,C/0,D/45,E/180,F/180,G/0,H/45}{\draw[fill](\i) circle (1pt)
						($(\i)+(\g:3mm)$)  node[scale=1]{$\i$};}
				\end{tikzpicture}
		\end{center}
		Ta có $A(0;0;0)$, $\overrightarrow{AE}=(0;0;1)$, $\overrightarrow{CD}=(-1;0;0)$.\\ 
		Đặt  $M(a;b;c)$, suy ra $\overrightarrow{AM}=(a;b;c)$.\\
		Đề cho $\overrightarrow{AM}+\overrightarrow{AE}=3\overrightarrow{CD}$,  ta được $\heva{&a+0=-3\\&b+0=0\\&c+1=0}\Rightarrow \heva{&a=-3\\&b=0\\&c=-1.}$\\
		 Suy ra $M(-3;0;-1)$.\\
		Mặt phẳng $(EBD)$, phương trình theo đoạn chắn là $x+y+z-1=0$.\\ 
		Khoảng cách từ điểm $M$  đến mặt phẳng $(EBD)$ bằng
		 \[\mathrm{d}\Big(M;(EBD)\Big)=\dfrac{\left|Ax+By+Cz+D\right|}{\sqrt{A^2+B^2+C^2}}=\dfrac{\left|-3+0-1-1\right|}{\sqrt{1^2+1^2+1^2}}=\dfrac{5\sqrt{3}}{3}.\]
		 Vậy khoảng cách bằng $\dfrac{5\sqrt{3}}{3}$.
	} 
\end{ex}

% nguồn: Phan2_Lop_12_Hinh_hoc_Chuong_5
\begin{bt}%[2H5V1-5]%[ma:48e31dc6]
	Trong không  gian với hệ trục $Oxyz$, khoảng cách giữa hai mặt phẳng $(P)\colon x+2y+2z-10=0$ và $(Q)\colon x+2y+2z-5=0$ bằng bao nhiêu?
	 \loigiai{
		\textbf{Cách 1}: Ta có $M(10; 0; 0) \in(P)$.  Vì $\dfrac{1}{1}=\dfrac{2}{2}=\dfrac{2}{2} \neq \dfrac{-10}{-5}$ nên hai mặt phẳng $(P)$ và $(Q)$ song song.\\
		Khi  đó, $d((P),(Q))=d(M,(Q))=\dfrac{|10-5|}{\sqrt{9}}=\dfrac{5}{3}$.\\
		\textbf{Cách 2}:  Công thức khoảng cách giữa hai mặt phẳng song song $(P)\colon a x+b y+c z-d_1=0$ và $(Q)\colon a x+b y+c z-d_2=0$ bằng
		\[\mathrm{d}((P),(Q))=\dfrac{\left|d_1-d_2\right|}{\sqrt{a^2+b^2+c^2}} \Rightarrow  \mathrm{d}((P),(Q))=\dfrac{|10-5|}{\sqrt{9}}=\dfrac{5}{3}. \]
	} 
\end{bt}

% nguồn: Phan2_Lop_12_Hinh_hoc_Chuong_5
\begin{bt}%[2H5V1-5]%[ma:92e9e340]
	Cho hình lập phương $ABCD.EFGH$ cạnh $1$. Điểm $M$ được cho thỏa mãn hệ thức véc-tơ $\overrightarrow{AM}+\overrightarrow{AE}=3\overrightarrow{CD}$. Tính khoảng cách  từ điểm $M$ đến mặt phẳng $(EBD)$.
	\loigiai{ 
		\begin{center}
			\begin{tikzpicture}[line join=round,line cap=round,>=stealth,scale=0.8]
					\coordinate (A) at (0,0); 
					\coordinate (y) at ($(A)+(0:5)$); 
					\coordinate (z) at  ($(A)+(90:5)$);
					\coordinate (x) at ($(A)+(-120:3)$); 
					\coordinate (B) at  ($(A)+(-120:1.5)$);
					\coordinate (D) at  ($(A)+(0:4)$);
					\coordinate (C) at  ($(B)+(D)-(A)$);
					\coordinate (E) at ($(A)+(90:4)$); 
					\coordinate  (F) at ($(B)+(90:4)$);
					 \coordinate (G) at ($(C)+(90:4)$);
					\coordinate  (H) at ($(D)+(90:4)$);
					 \fill[gray!20](B)--(D)--(E);
					\draw[->] (B) -- (x)  node[left] {$x$};
					\draw[->] (D) -- (y) node[below] {$y$}; 
					\draw[->] (E) -- (z) node[left]  {$z$};
					\draw (B)--(C)--(D) 
						(E)--(F)--(G)--(H)--cycle  (B)--(F)
						(C)--(G) 
						 (D)--(H);
					\draw  [dashed](A)--(B)
						 (A)--(D)
						 (A)--(E)
						 (B)--(D)--(E)--cycle;
					\foreach  \i/\g in {A/45,B/180,C/0,D/45,E/180,F/180,G/0,H/45}{\draw (\i) circle (1pt)
						 ($(\i)+(\g:3mm)$) node[scale=1]{$\i$};}
				\end{tikzpicture}
		\end{center}
		Ta  có $A(0;0;0)$, $\overrightarrow{AE}=(0;0;1)$, $\overrightarrow{CD}=(-1;0;0)$.\\
		Đặt $M(a;b;c)$,  suy ra $\overrightarrow{AM}=(a;b;c)$.\\
		Đề cho $\overrightarrow{AM}+\overrightarrow{AE}=3\overrightarrow{CD}$, ta  được $\heva{&a+0=-3\\&b+0=0\\&c+1=0}\Rightarrow \heva{&a=-3\\&b=0\\&c=-1.}$\\
		 Suy ra $M(-3;0;-1)$.\\
		Mặt phẳng $(EBD)$, phương trình  theo đoạn chắn là $x+y+z-1=0$.\\
		Khoảng cách từ điểm $M$ đến mặt phẳng  $(EBD)$ bằng
		\[\mathrm{d}\Big(M;(EBD)\Big)=\dfrac{\left|Ax+By+Cz+D\right|}{\sqrt{A^2+B^2+C^2}}=\dfrac{\left|-3+0-1-1\right|}{\sqrt{1^2+1^2+1^2}}=\dfrac{5\sqrt{3}}{3}.\] 
		Vậy khoảng cách  bằng $\dfrac{5\sqrt{3}}{3}$.
	} 
\end{bt}

% nguồn: Phan2_Lop_12_Hinh_hoc_Chuong_5
\begin{bt}%[2H5V1-5]%[ma:72c203b0]
	Cho hình lăng trụ tam giác đều $ABC.DEF$ biết rằng $AB=6$, $AD=2$. Gọi ba điểm $M$, $N$, $P$ lần lượt là trung điểm $DE$, $DF$,  $BC$. Lập hệ trục tọa độ $Oxyz$ như hình bên. Gọi điểm $S$ thỏa mãn hệ thức $\overrightarrow{SA}+2\overrightarrow{SB}+\overrightarrow{SC}=\overrightarrow{0}$. Tính khoảng cách từ điểm $S$ đến mặt phẳng $(MNP)$.
	 \loigiai{
		Ta có $PA=\dfrac{AB\sqrt 3}{2}=3\sqrt 3$; $PB=PC=\dfrac{BC}{2}=3$; $BE=BF=AD=2$.\\ 
		Đặt hệ trục tọa  độ như hình vẽ
		\begin{center}
			\begin{tikzpicture}[line join=round,line cap=round,>=stealth,scale=0.8]
					\coordinate (P) at  (0,0);
					\coordinate (y)  at ($(P)+(-5:4)$);
					\coordinate (z) at  ($(P)+(90:5)$);
					 \coordinate (x) at ($(P)+(-150:3)$);
					\coordinate (A) at ($(P)+(-5:3)$); 
					\coordinate  (B) at ($(P)+(-150:2)$);
					\coordinate  (C) at ($(P)+(30:2)$);
					\coordinate (D) at ($(A)+(90:3)$); 
					\coordinate (E) at  ($(B)+(90:3)$);
					\coordinate (F) at  ($(C)+(90:3)$);
					\coordinate (M) at ($(D)!0.5!(E)$); 
					\coordinate  (N) at ($(D)!0.5!(F)$);
					 \coordinate (P') at ($(E)!0.5!(F)$);
					\draw[->] (B) -- (x) node[left]  {$x$};
					\draw[->]  (A) -- (y) node[below] {$y$};
					\draw[->] (P') -- (z) node[left] {$z$}; 
					 \draw (D)--(E)--(F)--cycle (B)--(E)
						(A)--(D) 
						 (B)--(A)
						 (M)--(N);
					\draw [dashed] (B)--(C)--(A) 
						 (C)--(F)
						 (P)--(A)
						 (P)--(P')
						(M)--(P)--(N); 
					\foreach \i/\g in {A/-90,B/-90,C/30,D/0,E/180,F/90,M/90,N/30,P/180}{\draw (\i)  circle (1pt)
						 ($(\i)+(\g:3mm)$) node[scale=1]{$\i$};}
				\end{tikzpicture}
		\end{center}
		Khi đó $P(0;0;0)$, $A(3\sqrt{3};0;0)$, $B(0;3;0)$, $C(0;-3;0)$, $E(0;3;2)$, $F(0;-3;2)$, $D(3\sqrt{3};0;2)$.\\ 
		 Suy ra $M\left(\dfrac{3\sqrt{3}}{2};\dfrac{3}{2};2\right)$, $N\left(\dfrac{3\sqrt{3}}{2};\dfrac{-3}{2};2\right)$.\\
		Ta có $\overrightarrow{PM}=\left(\dfrac{3\sqrt{3}}{2};\dfrac{3}{2};2\right)$; $\overrightarrow{PN}=\left(\dfrac{3\sqrt{3}}{2};\dfrac{-3}{2};2\right)$.\\ 
		Suy ra $\left[\overrightarrow{PM},\overrightarrow{PN}\right]=\left(6;0;\dfrac{-9\sqrt{3}}{2}\right)\Rightarrow  \overrightarrow{n}=(4;0;-3\sqrt{3})$.\\
		Phương trình mặt phẳng $(MNP)$ là  $4x-3\sqrt{3}z=0$.\\
		Gọi $I$, $J$ lần lượt  là trung điểm cạnh $AC$, $IB\Rightarrow I\left(\dfrac{3\sqrt{3}}{2};\dfrac{-3}{2};0\right)$, $J\left(\dfrac{3\sqrt{3}}{4};\dfrac{3}{4};0\right)$.\\
		Ta có $\overrightarrow{SA}+2\overrightarrow{SB}+\overrightarrow{SC}=\overrightarrow{0}\Leftrightarrow  2\overrightarrow{SI}+2\overrightarrow{SB}=\overrightarrow{0}\Leftrightarrow 4\overrightarrow{SJ}=\overrightarrow{0}\Leftrightarrow S\equiv J$.\\
		Do đó  $S\left(\dfrac{3\sqrt{3}}{4};\dfrac{3}{4};0\right)$.\\
		Vậy khoảng cách từ điểm $S$ đến mặt phẳng $(MNP)$ là 
		\[\mathrm{d}\Big(S,(MNP)\Big)=\dfrac{\left|4\cdot \dfrac{3\sqrt 3}{4}-3\sqrt 3\cdot 0\right|}{\sqrt{4^2+\left(-3\sqrt 3\right)^2}}=\dfrac{3\sqrt{129}}{43}.\] 
	} 
\end{bt}

% nguồn: Phan2_Lop_12_Hinh_hoc_Chuong_5
\begin{bt}%[2H5C1-5]%[ma:ed1daec3]
	Trong không gian $Oxyz$, cho mặt phẳng $(P)\colon x+y+4z-3=0$ và điểm $A(1; 1; 3)$. Mặt phẳng $(Q) \parallel(P)$ và cắt các tia $Ox, Oy$ lần lượt tại các điểm $B$ và $C$ sao cho tam giác $ABC$ có diện tích bằng $2\sqrt{22}$. Khoảng cách từ điểm $M(2; 2; 1)$ đến $(Q)$ bằng bao  nhiêu?
	 \loigiai{
		Mặt phẳng $(Q) \parallel(P)  \Rightarrow(Q)$ có dạng: $x+y+4z+d=0\ (d \neq-3)$.\\
		Ta có $\begin{aligned}
		&(Q) \cap Ox=B(-d; 0; 0);\\ 
		&(Q) \cap Oy=C(0;-d; 0). 
	\end{aligned}$\\
	Do $B$, $C$ lần lượt thuộc các tia $Ox$, $Oy$ nên suy ra  $d < 0$.\\
	Ta có  $\overrightarrow{AB}=(-d-1;-1;-3)$, $\overrightarrow{AC}=(-1;-d-1;-3) \Rightarrow [\overrightarrow{AB}, \overrightarrow{AC}]=\left(-3d;-3d; d^2+2d\right)$.
	\allowdisplaybreaks 
	\begin{eqnarray*}
		S_{\triangle ABC}=2\sqrt{22} &\Leftrightarrow &\dfrac{1}{2}\left|[\overrightarrow{AB},  \overrightarrow{AC}]\right|=2\sqrt{22}\\
		&\Leftrightarrow &9d^2+9d^2+\left(d^2+2d\right)^2=352\\ 
		&\Leftrightarrow  &d^4+4d^3+22d^2-352=0\qquad (*)
	\end{eqnarray*}
	Giải  (*) chỉ có $d=-4$ thỏa mãn.\\
	Khi đó ta có phương trình mặt phẳng $(Q)\colon  x+y+4z+4=0$.\\
	Khoảng cách từ  điểm $M(2; 2; 1)$ đến $(Q)$ bằng $\dfrac{|2+2+4\cdot 1+4|}{\sqrt{1^2+1^2+4^2}}=2\sqrt{2}$.
 }
\end{bt}
